% Anonymous edition. Compile both documents with python3 build.py.
% Official econsocart class/configuration are redistributed without modification.
\documentclass[ecta,draft,supplement]{econsocart}
\ifdefined\pdfminorversion\pdfminorversion=7\fi
\usepackage[utf8]{inputenc}
\usepackage{microtype,mathtools,booktabs,array,enumitem,needspace,longtable,tabularx}
\usepackage[hidelinks,hypertexnames=false]{hyperref}
\usepackage{xurl,bookmark}
% Portable mappings for font shapes used by the official public class.
\DeclareFontShape{T1}{ptm}{m}{uc}{<->ssub*ptm/m/n}{}
\DeclareFontShape{T1}{ptm}{m}{up}{<->ssub*ptm/m/n}{}
\DeclareFontShape{T1}{ptm}{m}{scit}{<->ssub*ptm/m/sc}{}
\hypersetup{pdftitle={The Cost of Voluntary Exchange: Allocation and Settlement under Separate Public Accounts},pdfauthor={},pdfsubject={Anonymous manuscript},pdfkeywords={voluntary exchange; separate public accounts; settlement; reserves}}
\setlength{\emergencystretch}{2em}
\clubpenalty=10000
\widowpenalty=10000
\displaywidowpenalty=10000
\setlist{itemsep=3pt,topsep=5pt}
\setcounter{secnumdepth}{2}
% Neutral masthead for an anonymous working manuscript. No submission is asserted.
\makeatletter
\def\copyright@text{Anonymous Manuscript}
\makeatother
\startlocaldefs
\newtheorem{theorem}{Theorem}
\newtheorem{proposition}[theorem]{Proposition}
\newtheorem{lemma}[theorem]{Lemma}
\newtheorem{corollary}[theorem]{Corollary}
\theoremstyle{definition}
\newtheorem{definition}[theorem]{Definition}
\theoremstyle{definition}
\newtheorem{remark}[theorem]{Remark}
\newtheorem{example}[theorem]{Example}
\newcommand{\R}{\mathbb R}
\newcommand{\FS}{\mathcal F^{S}}
\newcommand{\FP}{\mathcal F^{P}}
\newcommand{\FB}{\mathcal F^{B}}
\newcommand{\argmax}{\operatorname*{arg\,max}}
\newcommand{\argmin}{\operatorname*{arg\,min}}
\newcommand{\supp}{\operatorname{supp}}
\renewcommand{\thetheorem}{A.\arabic{theorem}}

\endlocaldefs
\begin{document}
\begin{frontmatter}
\title{Online Appendix: The Cost of Voluntary Exchange:\\
Allocation and Settlement under Separate Public Accounts}
\runtitle{Online Appendix: The Cost of Voluntary Exchange}
\end{frontmatter}
\noindent Supplementary proofs and extensions for the accompanying main manuscript. Cross-document references refer to the revised files.
\begin{appendix}
\numberwithin{equation}{section}
\section{General reserve bounds and exact recovery}
\label{app:designproofs}

Both results retain $B\equiv0$, complete information, positive finite
action caps, linear resource effects, divisible irreversible adjustments,
complete bilateral access, demander-only transactions, and enforceable
compensation balanced within each approval. The second result additionally
requires full reserve repayment.

\subsection{Bounds on the approval count}

Initially let employee benefits and costs be linear, with $v_i,c_j>0$,
and put $\bar v=\max_i v_i$. A pair increment $(x,y)$ is voluntary
exactly when $v_ix\geq c_jy$. Let $R^0\geq0$,
$Z=\{k:R_k^0=0\}$, and $R^\epsilon=R^0+\epsilon r$, where
$r\geq0$ and $r_k>0$ for every $k\in Z$. Assume a blocked start:
no nonzero voluntary approval is feasible at $R^0$, although
$V^S(R^0)>0$. Define the compact normalized direction set
\[
 \mathcal D=\{e_i:i\in D\}\cup
 \bigcup_{i,j}\{xe_i+ye_j:x,y\geq0,\ x+y=1,\ v_ix\geq c_jy\}.
\]
For any $\lambda\gg0$ on the initially empty accounts, put
\begin{equation}
 \delta_\lambda=\min_{d\in\mathcal D}
       \sum_{k\in Z}\lambda_k[-(Md)_k]_+,\qquad
 g_\lambda=\max\{0,\max_{d\in\mathcal D}\lambda^{\mathsf T}M_Zd\},
 \qquad S_\lambda=\lambda^{\mathsf T}r_Z.
 \label{eq:generalseedconstants}
\end{equation}
These analytical weights do not permit transfers between accounts.

\begin{theorem}[General approval bounds]
\label{thm:generalseedbounds}
Under these assumptions, $\delta_\lambda,S_\lambda>0$, and
for every $\epsilon>0$ and integer $T\geq0$,
\begin{equation}
 V_T^{\mathrm{seq}}(R^\epsilon)\leq
 \begin{cases}
 \displaystyle\frac{\bar v\epsilon S_\lambda}{g_\lambda}
       [(1+g_\lambda/\delta_\lambda)^T-1],&g_\lambda>0,\\[5pt]
 \displaystyle\frac{\bar v\epsilon S_\lambda}{\delta_\lambda}T,
       &g_\lambda=0.
 \end{cases}
 \label{eq:generalseedlower}
\end{equation}
Fix a planner optimum $z^*$ at $R^0$ and
$0<\eta<V^S(R^0)$. The target
$z^\eta=qz^*$, where $q=\eta/V^S(R^0)$, admits exact voluntary
implementation in at most
\begin{equation}
 B_\eta/\epsilon+C_\eta
 \label{eq:constructiveseedupper}
\end{equation}
approvals, for finite computable constants independent of $\epsilon$.
If
\begin{equation}
 \lambda^{\mathsf T}M_Zd\leq0\quad(d\in\mathcal D),
 \label{eq:nonexpandingcertificate}
\end{equation}
then, writing $T_\eta$ for the minimum approvals achieving surplus
at least $\eta$ and $N_\eta$ for those implementing $z^\eta$,
\begin{equation}
 \frac{\eta\delta_\lambda}{\bar vS_\lambda\epsilon}
 \leq T_\eta(\epsilon)\leq N_\eta(\epsilon)
 \leq\frac{B_\eta}{\epsilon}+C_\eta.
 \label{eq:sharpinverseseed}
\end{equation}
Thus both requirements have inverse-reserve order under the certificate.
\end{theorem}

\begin{proof}
If $\delta_\lambda=0$, compactness gives $d\in\mathcal D$ with
$M_Zd\geq0$. A sufficiently small positive multiple respects every
positive initial account and action cap, contradicting the blocked
start. That assumption also implies $Z\ne\varnothing$, so
$S_\lambda>0$. Write an approval as $s_td_t$, $d_t\in\mathcal D$,
and let $H_t=\lambda^{\mathsf T}(R^\epsilon+Mz_t)_Z$.
Accountwise feasibility implies
\[
 \delta_\lambda s_t\leq H_{t-1},\qquad
 H_t\leq H_{t-1}+g_\lambda s_t,
 \qquad H_0=\epsilon S_\lambda.
\]
Iterating and summing these inequalities, then using
$W(z_T)\leq\bar v\sum_{t\leq T}s_t$, proves
\eqref{eq:generalseedlower}.

For the upper bound, use the decomposition proved below in
Theorem~\ref{thm:temporaryreserve}. Its marginal-financing condition
holds for the linear target because $Q-C=W(z^\eta)=\eta>0$.
This gives voluntary bundles $u^1,\ldots,u^L$ summing to $z^\eta$,
where $L\leq|D||S|+|D|$. Set
\[
 A_k=\sum_\ell[-(Mu^\ell)_k]_+,\qquad
 a_\eta=\max_{k\in Z}A_k/r_k,\qquad
 b_\eta=\max_{k\notin Z}\frac{A_k}{(1-q)R_k^0},
\]
with an empty maximum equal to zero, and choose
\begin{equation}
 n=\max\left\{1,\left\lceil
       \max_k\frac{A_k}{(1-q)R_k^0+\epsilon r_k}
       \right\rceil\right\}.
 \label{eq:explicitseedrounds}
\end{equation}
All denominators are positive. Execute $n$ identical ordered rounds
of $u^1/n,\ldots,u^L/n$. A round starting at fraction $t$ of
$z^\eta$ has balances
$(1-tq)R^0+tq(R^0+Mz^*)+\epsilon r\geq(1-q)R^0+\epsilon r$.
Each within-round decrease is at most $A_k/n$, so every prefix is
feasible. Activities increase to the target and every divided bundle
remains voluntary. The rounds follow
\citet[Lemma~13]{FracaHaddad2015}; the decomposition verifies compensation.

Since $n\leq a_\eta/\epsilon+1+b_\eta$, take
$B_\eta=La_\eta$ and $C_\eta=L(1+b_\eta)$. Finally,
\eqref{eq:nonexpandingcertificate} implies $g_\lambda=0$;
inserting the required surplus $\eta$ in
\eqref{eq:generalseedlower} gives \eqref{eq:sharpinverseseed}.
\end{proof}

\subsection{Exact recovery with repayment}

\begin{theorem}[Exact recovery with a repayable reserve]
\label{thm:temporaryreserve}
Retain the section's institutional assumptions. Let $R^0\geq0$
and $r\gg0$. Benefits are differentiable, concave, and
nondecreasing; costs are differentiable and convex with strictly
positive derivatives, with $U_i(0)=C_j(0)=0$. Fix any original-ledger
planner optimum $z^*=(a^*,b^*)$ with
$V_0=W(z^*)=V^S(R^0)>0$.
For every $\epsilon>0$, a finite sequence of voluntary approvals
with balanced employee compensation implements $z^*$ from
$R^0+\epsilon r$ and returns the full reserve. Hence the original
planner value is exactly recoverable before implementation costs.

There is a decomposition $z^*=\sum_{\ell=1}^Lu^\ell$ into
nonzero bilateral or demander-only bundles, with
$L\leq|D||S|+|D|$. Define
$A_k^*=\sum_\ell[-(Mu^\ell)_k]_+$ and
$H^*=\max_k A_k^*/r_k$. The minimum approvals implementing this
optimum with repayment satisfy
\[
 N^*(\epsilon)\leq L\max\{1,\lceil H^*/\epsilon\rceil\}
                  \leq LH^*/\epsilon+L.
\]
With additionally positive linear preferences, a blocked start,
and \eqref{eq:nonexpandingcertificate}, also
$N^*(\epsilon)\geq V_0\delta_\lambda/(\bar vS_\lambda\epsilon)$,
using \eqref{eq:generalseedconstants}. The inverse-reserve order is
therefore necessary and sufficient in that class.
\end{theorem}

\begin{proof}
Let $g_i=U_i'(a_i^*)\geq0$, $h_j=C_j'(b_j^*)>0$,
$Q=\sum_i g_ia_i^*$ and $C=\sum_jh_jb_j^*$.
The segment from $0$ to $z^*$ is feasible. Optimality at its
endpoint gives the inward radial condition $Q\geq C$.
If $C>0$, also $Q>0$, and define pair increments
\[
 x_{ij}=a_i^*h_jb_j^*/C,\qquad
 y_{ij}=b_j^*g_ia_i^*/Q
\]
for each demander and active supplier. They satisfy
$\sum_jx_{ij}=a_i^*$, $\sum_i y_{ij}=b_j^*$, and
$g_ix_{ij}\geq h_jy_{ij}$. When $g_i=0$ the bundle is
demander-only. If $C=0$, strictly positive $h_j$ implies
$b^*=0$, and one demander-only bundle per active demander suffices.
Discarding zero bundles gives the stated $u^\ell$ and bound on $L$.

Execute $n=\max\{1,\lceil H^*/\epsilon\rceil\}$ rounds
of $u^1/n,\ldots,u^L/n$. At fraction $t$ of completed rounds,
balances are $(1-t)R^0+t(R^0+Mz^*)+\epsilon r\geq\epsilon r$.
The greatest within-round decrease is at most
$A_k^*/n\leq\epsilon r_k$ in each account. All intermediate
balances are nonnegative, and activities increase coordinatewise
from zero to $z^*$, respecting caps. Linear resource effects
preserve feasibility within each approval.

At any such cumulative allocation, concavity and convexity imply
\[
 \Delta U_i\geq g_ix_{ij}/n
             \geq h_jy_{ij}/n\geq\Delta C_j.
\]
Charge the demander exactly the supplier's incremental cost and
pay it to that supplier. Compensation balances and both weakly gain.
Demander-only increments have nonnegative benefits and need no
transfer. Thus participation holds at each approval under nonlinear
preferences.

The final balances are $R^0+Mz^*+\epsilon r$, so the reserve
is fully repayable. Conversely, repayment at any terminal allocation
$z$ implies $R^0+Mz\geq0$ and therefore $W(z)\leq V^S(R^0)$.
There are $Ln\leq LH^*/\epsilon+L$ approvals. Under the additional
linear, blocked-start, and nonexpansion assumptions,
\eqref{eq:generalseedlower} with surplus $V_0$ and
$g_\lambda=0$ supplies the positive matching lower bound.
\end{proof}

\section{Exact reserve and cost frontiers}\label{app:localrobustness}
\subsection{Proof of the exact productive frontier}

\begin{proof}[Proof of Proposition~2%<xref:prop:exactproductive>
]
First take $h>0$ and abbreviate $\gamma=\gamma_h$ and
$\epsilon=e$. At any cumulative allocation, write $P,W,H$ for the
three balances and set
\[
 S=P+W+H=3e+hb_1,\qquad
 q=W-H=\frac{b_1-b_2}{2},
\]
\[
 u=S-q=P+2H,\quad v=S+q=P+2W.
\]
Feasibility implies $u,v\geq0$, with initial state $(3e,3e)$.
A supplier-1 increment $y$ changes $(u,v)$ by
$(-(1/2-h)y,(1/2+h)y)$: it converts some of $u$ into $v$ at rate
$\gamma$. A supplier-2 increment changes $(u,v)$ by
$(y/2,-y/2)$ and converts some of $v$ into $u$ at unit rate.
Demander-only transactions leave both coordinates unchanged. The target
requires
\[
 u=v=s,\qquad s=3e+h\theta/2.
\]

These two-coordinate restrictions give a lower bound for every physical
schedule. In this relaxed problem, delete transactions with no supplier
increment and combine adjacent increments of the same supplier. This
deletion is not asserted to preserve the original physical balances.
The remaining word alternates and has length at most $T$; it has at least
two steps because the target uses both suppliers. To bound its final
common stock, work backward.
If its last transaction uses supplier 1, equality of the terminal
coordinates gives $s=(\gamma u+v)/(1+\gamma)$ at the preceding state.
If it uses supplier 2, it gives $s=(u+v)/2$. For a nonnegative linear
bound $A u+B v$, a preceding supplier-1 conversion changes the coefficient
on $u$ to $\max\{A,\gamma B\}$; a preceding supplier-2 conversion
changes the coefficient on $v$ to $\max\{B,A\}$. These replacements
follow by maximizing the converted amount over $[0,u]$ and $[0,v]$,
respectively, so they bound every admissible partial conversion.

For an alternating word, the replacements select full conversion at every
prefinal step. Evaluated at $(3e,3e)$, a word beginning with supplier 1
and of length $2n$ has bound
$3e(1+\gamma)\gamma^{n-1}/2$; one beginning with supplier 2 has
bound $3e\,2\gamma^n/(1+\gamma)$. The first is weakly larger,
because $(1+\gamma)^2\geq4\gamma$. At length $2n+1$, both bounds
are $3e\gamma^n$. Thus $s\leq3eG_T(h)$ for length $T$.
The sequence $G_T(h)$ is strictly increasing, so the same bound covers
all shorter words. Substituting the required $s$ gives
$e\geq h\theta/[6(G_T(h)-1)]$.

It remains to lift an attaining conversion path to physical, voluntary
transactions. Choose $e$ equal to this lower bound, begin with supplier 1,
and alternate suppliers. Make every prefinal conversion full. For a
prefinal supplier-1 step, select the demander increment so that its endpoint
has $P=H=0$; equivalently, its new cumulative demander quantity is
$a=b_1+b_2-e$. For a prefinal supplier-2 step, select the increment so
that $P=W=0$; equivalently, set $a=b_1/2+3b_2/2-e$.
The other account is nonnegative in each case. The first step has
$y=6e/(1-2h)$ and demander increment $x=y-e>0$.
An interior supplier-1 step has $x=(3+2h)y/4$, and an interior
supplier-2 step has $x=5y/4$.

Use the last conversion to equalize $u$ and $v$. The amplification
calculation gives $u=v=3eG_T(h)=3e+h\theta/2$. The identities for
$S$ and $q$ therefore give $b_1=b_2=\theta/2$.
Set the final demander quantity to $a=\theta$. If the last supplier
is 1, its increments satisfy $x=y/2+e$; if it is 2, they satisfy
$x=y+e$. The resulting final balances are
$(e+h\theta/2,e,e)$, so the full reserve can be returned.

Every displayed increment is nonnegative and satisfies $x>y/10$.
Paying the supplier $y/10$ and charging the demander the same amount
makes each transaction weakly voluntary and exactly balanced. Since all
coordinates increase to the target, every prefix respects the unit action
caps. The prefinal endpoints and the final endpoint are resource-feasible;
linear account effects also preserve feasibility along each divisible
transaction. Any larger reserve permits the same schedule with the excess
capacity retained throughout.

At $h=0$, the conserving-economy proof gives
$e=\theta/[12(T-1)]$. Finally,
$\gamma_h=1+4h+O(h^2)$ and
$G_T(h)=1+2h(T-1)+O(h^2)$ for fixed $T$, which gives the asserted
continuous limit.
\end{proof}

\subsection{The conserving frontier}
\label{app:exactseed}

\begin{proof}[Proof of \eqref{eq:exactseedthreshold}]
Fix $\theta\in(0,1]$ and a schedule with at most $T$ bilateral
settlements that implements $z_\theta=(\theta,\theta/2,\theta/2)$
from $\epsilon(1,1,1)$. Zero increments can be appended to give exactly
$T$ steps. Since both suppliers must make positive adjustments, $T\geq2$.
Write the resource balances after step $t$ as $(P_t,W_t,H_t)$ and put
$q_t=W_t-H_t$. Every resource column sums to zero, so feasibility gives
\[
 P_t+W_t+H_t=3\epsilon,
 \qquad -3\epsilon\leq q_t\leq3\epsilon.
\]
The target has zero effect on every resource account; hence
$q_0=q_T=0$. A supplier-1 transaction of size $(x,y)$ changes $q_t$
by $y/2$, while a supplier-2 transaction changes it by $-y/2$.
The demander's increment cancels from $W_t-H_t$. Therefore
\[
 \theta=b_1+b_2
   =2\sum_{t=1}^T|q_t-q_{t-1}|
   \leq 2\{3\epsilon+6\epsilon(T-2)+3\epsilon\}
   =12\epsilon(T-1).
\]
This proves the lower bound for every voluntary schedule, regardless of
the supplier order or the division of the demander's adjustment.

For attainment, set $e=\theta/[12(T-1)]$. Start with a supplier-1
transaction $(x,y)=(5e,6e)$. It changes the resource balance from
$(e,e,e)$ to $(0,3e,0)$. Each interior transaction alternates suppliers:
supplier 2 uses $(15e,12e)$ and changes $(0,3e,0)$ to $(0,0,3e)$;
supplier 1 uses $(9e,12e)$ and reverses that change. After $T-2$
interior transactions, finish with supplier 1 and $(4e,6e)$ if $T$
is odd, or with supplier 2 and $(7e,6e)$ if $T$ is even. Either final
transaction returns the balance to $(e,e,e)$.

The first and last increments of $q_t$ have absolute size $3e$,
and every interior increment has absolute size $6e$. Hence
$b_1+b_2=12e(T-1)=\theta$. Since $q_T=q_0$, the two suppliers'
totals are equal, giving $b_1=b_2=\theta/2$. The unchanged final
position balance then gives $a=b_1+b_2=\theta$.
All increments are nonnegative, so the target respects every individual
cap and so does each prefix. Each displayed transaction satisfies
$x-y/10>0$ and therefore admits balanced compensation with both employees
weakly better off. Every prefix has nonnegative public balances. Thus the
lower bound is attained for every $\theta\in(0,1]$ and integer $T\geq2$.
Any larger reserve also permits the same schedule, completing the proof.
\end{proof}

\subsection{Exact costs and the fixed-target comparison}
\begin{proof}[Proof of Theorem~11%<xref:thm:productivecost>
]
An admissible $T$-approval schedule requires reserve at least
$\epsilon_T$, and Proposition~2%<xref:prop:exactproductive>
 attains this amount.
Taking infima proves \textup{(5)}%<xeqref:eq:productivecostexact>
. The minimum exists because
the cost is at least $\tau T$. A supplier-1 pair $(\theta/2,\theta/2)$ followed
by the same supplier-2 pair is voluntary and feasible from $\theta/4$ in each
account, and returns the reserve. No schedule uses fewer than two approvals,
so any larger reserve is dominated.

Under conservation, writing $n=T-1$ gives the scalar objective
$\rho\theta/(12n)+\tau(n+1)$; the arithmetic--geometric mean inequality and
rounding its real minimizer establish the stated order. For fixed $h>0$,
the formula for $G_T$ gives constants $A_1,A_2,b>0$ with
$A_1\exp(-bT)\leq\epsilon_T\leq A_2\exp(-bT)$ for every $T\geq2$.
Choose $T=\lceil b^{-1}\log(\rho/\tau)\rceil+2$ for an upper bound.
For the lower bound, if $T\geq(2b)^{-1}\log(\rho/\tau)$ the approval cost
already has the required order. Otherwise the reserve cost is at least
$A_1\sqrt{\rho\tau}$, which exceeds a constant multiple of
$\tau\log(\rho/\tau)$ as $\tau/\rho\downarrow0$. This proves both bounds.
\end{proof}

\begin{corollary}[A strict reversal of the preferred settlement rule]
\label{cor:productivereversal}
Fix $\theta=1$, $\rho=12$, $\tau=1/10000$, $F=1/100$, and $\mathcal K=0$.
At $h=0$, the sequential minimum is attained at reserve $1/1200$ with
$101$ approvals, and
\[
 Q_0=\frac{201}{10000}
 >\frac{101}{10000}=\tau+F.
\]
At $h=1/4$, the sequential minimum is attained at reserve $1/104952$
with $16$ approvals, and
\[
 Q_{1/4}=\frac{9371}{5466250}
 <\frac{101}{10000}=\tau+F.
\]
Joint settlement is therefore strictly cheaper in the conserving economy,
while sequential settlement is strictly cheaper after the productive
perturbation. In each economy, implementing the target is strictly better
than retaining the baseline under the stated cost criterion.
\end{corollary}

\begin{proof}[Proof of Corollary~\ref{cor:productivereversal}]
The exact conserving threshold \eqref{eq:exactseedthreshold} gives
\[
 Q_0=\tau+\min_{n\geq1}
      \left\{\frac{\rho}{12n}+\tau n\right\}.
\]
For the specified prices, the expression in braces is
$1/n+n/10000$. Its unique real minimizer is $n=100$, which is an
integer. The minimum is $1/50$, giving $Q_0=201/10000$,
$101$ approvals, and reserve $1/1200$.

For $h=1/4$, the exact threshold in
\textup{(4)}%<xeqref:eq:exactproductiveseed>
 gives the cost for $T$ approvals as
\[
 c_T=\frac1{2(g_T-1)}+\frac{T}{10000},\qquad
 g_{2n}=2\cdot3^{n-1},\quad g_{2n+1}=3^n.
\]
Within either parity, increasing $T$ by two multiplies $g_T$ by three.
The corresponding cost difference is
\[
 \frac1{5000}-\frac{g_T}{(3g_T-1)(g_T-1)}.
\]
These differences strictly increase as $g_T$ increases, because
$g/[(3g-1)(g-1)]$ is strictly decreasing for $g\geq2$.
For the even subsequence the difference is negative at $g_{14}=1458$
and positive at $g_{16}=4374$, so its unique minimum is at $T=16$.
For the odd subsequence it is negative at $g_{13}=729$ and positive
at $g_{15}=2187$, so its unique minimum is at $T=15$.
Finally,
\[
 c_{15}-c_{16}=\frac{687811}{47796890000}>0.
\]
Hence $T=16$ is the global minimizer. Its threshold reserve is
$1/[24(4374-1)]=1/104952$, and its cost is
$1/8746+16/10000=9371/5466250$.
The common joint cost is $101/10000$, strictly between the two
exact sequential minima. Both preferred procedures cost less than
the target surplus $9/10$, and $\mathcal K=0$, establishing adoption as well.
\end{proof}

\subsection{Endogenous choice in the conserving economy}

The preceding theorem solves the conserving allocation choice at the stated
prices. The same argument yields a complete comparative-static solution for
$M_0$. Repayment requires $M_0z\geq0$; zero column sums force equality.
Thus every terminal allocation is $z_\theta$, $0\leq\theta\leq1$.
The frontier specializes to
\begin{equation}
 \epsilon_T(\theta)=\frac{\theta}{12(T-1)},\qquad T\geq2.
 \label{eq:exactseedthreshold}
\end{equation}
Allow common verification cost $\mathcal K(z_\theta)=k\theta$, $k\geq0$, and
retain the same fixed $F\geq0$ for every nonzero joint package. The baseline
incurs no cost. If a reserve cap is imposed, assume it is at least $1/12$,
so it does not restrict any minimum reserve in this comparison.

\begin{proposition}[Endogenous quantity, reserve, and approval rule]
\label{prop:exactdesign}
Under these conditions, define
\begin{equation}
 G(\rho,\tau)=\min_{n\in\{1,2,\ldots\}}
       \left\{\frac{\rho}{12n}+\tau n\right\}.
 \label{eq:exactG}
\end{equation}
The authority's maximum net value, allowing either institution and the
baseline, is
\begin{equation}
 \mathcal V(\rho,\tau,F,k)=
 \max\left\{0,\frac9{10}-k-\tau-
                         \min\{F,G(\rho,\tau)\}\right\}.
 \label{eq:exactdesignvalue}
\end{equation}
Whenever positive net surplus is attainable, an optimal allocation is
$z^*=(1,1/2,1/2)$. Joint approval implements it at lower cost when
$F<G(\rho,\tau)$, and sequential approval does so at lower cost when
$F>G(\rho,\tau)$. Under sequential approval, an integer $n^*\geq1$
minimizes \eqref{eq:exactG} if and only if
\begin{equation}
 \tau n^*(n^*-1)\leq\frac\rho{12}
       \leq\tau n^*(n^*+1).
 \label{eq:integerreservechoice}
\end{equation}
Conditional on implementing $z^*$ sequentially, the least-cost choice uses $n^*+1$ approvals and reserve
$\epsilon^*=1/(12n^*)$.
\end{proposition}


\begin{proof}[Proof of Proposition~\ref{prop:exactdesign}]
Fix $\theta>0$ and $T=n+1\geq2$. Since reserve is costly, its optimal
amount for that pair $(\theta,n)$ is the minimum in
\eqref{eq:exactseedthreshold}. Sequential net value is therefore
\[
 \theta\left(\frac9{10}-k-\frac\rho{12n}\right)-\tau(n+1).
\]
For a fixed $n$, if the coefficient of $\theta$ is positive, this expression
is maximized at $\theta=1$. If it is nonpositive, every positive allocation
gives negative net value and the baseline is better. Joint net value is
$\theta(9/10-k)-\tau-F$ and has the same full-quantity-or-baseline
property. Maximizing over $n$, the two institutions, and the baseline gives
\eqref{eq:exactdesignvalue}.

Write $f(n)=\rho/(12n)+\tau n$. Since
\[
 f(n+1)-f(n)=\tau-\frac\rho{12n(n+1)},
\]
its consecutive differences increase with $n$. A global minimum is
characterized by $f(n)-f(n-1)\leq0\leq f(n+1)-f(n)$ for $n\geq2$,
and just the second inequality for $n=1$. These conditions are exactly
\eqref{eq:integerreservechoice}. The minimum exists because
$f(n)\to\infty$ as $n\to\infty$.
\end{proof}

\section{Execution of productive releases}
\label{app:generalproductive}

\begin{proof}[Proof of Theorem~4%<xref:thm:generalproductive>
]
Use the objects in Definition~3%<xref:def:executableproductive>
 and set
$s_0=\sigma\epsilon$, so $\epsilon r\geq s_0p$.
After productive bundles of cumulative size $a\leq\alpha$, balances are
$\epsilon r+ap$. A further repetition $tu^1,\ldots,tu^k$ is feasible whenever
\[
 \begin{gathered}
 0<t\leq\min\{\alpha-a,(s_0+a)/\kappa\},\\
 \epsilon r+ap+tM\sum_{j=1}^s u^j
 \geq(s_0+a-t\kappa)p\geq0.
 \end{gathered}
\]
Its terminal balances are $\epsilon r+(a+t)p$. Choosing
$t=(s_0+a)/\kappa$ multiplies $s_0+a$ by $\Lambda=1+1/\kappa$.
Consequently, maximal repetitions followed by a truncated final repetition
reach $a=\alpha$ exactly in
\[
 n=\left\lceil
       \frac{\log(1+\alpha/(\sigma\epsilon))}{\log\Lambda}
     \right\rceil
\]
repetitions. The completion list then has prefix balances
$\epsilon r+\alpha p+M\sum_{j=1}^s v^j\geq\epsilon r$ and terminates at
$z^*=\alpha d+c$. Returning the reserve leaves $Mz^*\geq0$.
Thus the total number of approvals is at most $kn+\ell$.

All increments are nonnegative and sum to $z^*$, so every partial allocation
respects its caps. Linear interpolation preserves account feasibility within
each approval. For a bilateral increment $(x,y)$ involving demander $i$ and
supplier $j$, transfer $c_jy$ from the demander to the supplier: the gains
are $v_ix-c_jy\geq0$ and zero, respectively. Demander-only increments need
no transfer. Compensation is therefore voluntary and exactly balanced.

For the lower bound, fix any $\lambda\gg0$. Blocked start and compactness
give $\delta_\lambda>0$ in \eqref{eq:generalseedconstants}. Also
$g_\lambda>0$: otherwise $\lambda^{\mathsf T}Mu^j\leq0$ for every bundle,
contradicting $\lambda^{\mathsf T}\sum_jMu^j=\lambda^{\mathsf T}p>0$.
Theorem~\ref{thm:generalseedbounds} implies
\[
 T\geq
 \frac{\log\!\left(1+
 g_\lambda W(z^*)/(\bar v\epsilon S_\lambda)\right)}
 {\log(1+g_\lambda/\delta_\lambda)}
 =\Omega(\log(1/\epsilon)).
\]
This matches the upper order. With a nonexpansion certificate instead,
$g_\lambda=0$ and the same theorem gives
$T\geq W(z^*)\delta_\lambda/(\bar vS_\lambda\epsilon)$.
For a planner optimum, Theorem~\ref{thm:temporaryreserve} supplies a repayable
$O(1/\epsilon)$ schedule under the maintained complete bilateral access.
This proves the contrasting inverse-reserve order. The summation argument
above also proves incompatibility of the two certificates; it does not
make either one necessary for all technologies.
\end{proof}

\paragraph{Membership of \texorpdfstring{$M_h$}{Mh}.}
For $0<h<1/2$, choose
\[
 \begin{gathered}
 u^1=(1,1,0)^{\mathsf T},\quad u^2=(1,0,1)^{\mathsf T},\\
 d=(2,1,1)^{\mathsf T},\quad p=(h,0,0)^{\mathsf T},\quad
 \kappa=\frac{1/2-h}{h}.
 \end{gathered}
\]
The bundles each have surplus $9/10$, and their resource changes are
$(h-1/2,1/2,0)^{\mathsf T}$ and $(1/2,-1/2,0)^{\mathsf T}$.
Thus every prefix is at least $-\kappa p$. Take
\[
 \alpha=\frac1{2+3h},\qquad
 c=v^1=\frac{(3h,0,2h)^{\mathsf T}}{2+3h},\qquad
 M_hc=\frac{(-h,0,h)^{\mathsf T}}{2+3h}.
\]
The completion bundle has positive surplus $14h/[5(2+3h)]$, satisfies
$\alpha p+M_hc\geq0$, and gives $\alpha d+c=z_h^*$.
Only scaled bundles are executed; their nonnegative cumulative activities
are bounded by $z_h^*$ even though the auxiliary vector $d$ exceeds a unit cap.
For $r=(1,1,1)^{\mathsf T}$, $\sigma=1/h$ and $\Lambda=1/(1-2h)$.

The zero start is blocked throughout $0\leq h<1/2$. A supplier-1 pair would
require $x\geq y$ and $x\leq(3/4+h/2)y<y$; a supplier-2 pair would require
$x\geq3y/2$ and $x\leq5y/4$. A positive demander-only move consumes pension
capacity. At $h=0$, every column sums to zero, giving the nonexpansion
certificate $(1,1,1)^{\mathsf T}$. The general order comparison therefore
applies to $(M_0,M_h)$; the exact constants and parity terms in
Proposition~2%<xref:prop:exactproductive>
 require the separate frontier proof.

\section{A primitive LP certificate and an interacting release cycle}
\label{app:primitiveproductivity}

\begin{proof}[Proof of Proposition~5%<xref:prop:primitiveproductivity>
]
The planner LP attains $V^*$ on its compact activity box. For a
fixed support and order, the second program is also a linear program
with compact feasible set: nonnegative bundle activities sum to a
cap-bounded vector, and the nonempty support bounds $\eta$ through
$Mz$. Thus any positive value $L_{R,\pi}$ is attained.

At an attaining solution, feasibility and $W(z)=V^*$ make $z$ a
planner optimum. Its release $p=Mz$ is positive on $R$ and zero
elsewhere. Every bundle is voluntary. Outside $R$, the LP makes
every resource prefix nonnegative; inside $R$, the finite maximum
defining $\kappa$ bounds every prefix below by $-\kappa p$.
If $\kappa=0$, the first nonzero bundle would be feasible from
zero, contradicting the blocked start. Such a bundle exists since
$V^*>0$. Also $\sigma>0$ because $r\gg0$ and $p>0$ on $R$.
Delete zero bundles. The remaining list satisfies
Definition~3%<xref:def:executableproductive>
 with $d=z$, $\alpha=1$,
and no completion. Theorem~4%<xref:thm:generalproductive>
 gives the
stated count, matching lower bound, balanced compensation, and
repayment. Its cumulative activities stay below $z$, hence below
the caps. The assumed comparison technology then meets
Corollary~8%<xref:cor:optimizedreversal>
, proving the institutional claim.
\end{proof}

The finite enumeration searches over supports and orders. More general
certificates can use repeated transaction types or a separate completion.

\paragraph{The original family.}
For $0<h<1/2$, write $D_h=2+3h$ and choose
$u^1=D_h^{-1}(3/4+h/2,1,0)^{\mathsf T}$ and
$u^2=D_h^{-1}((5/4)(1+2h),0,1+2h)^{\mathsf T}$.
They sum to $z_h^*$ and are voluntary: their demander-to-supplier
ratios exceed $1/10$. Direct multiplication gives
\[
 M_hu^1=\frac{(0,1+2h,-1+2h)^{\mathsf T}}{4D_h},
 \qquad
 M_hu^2=\frac{(0,-1-2h,1+2h)^{\mathsf T}}{4D_h}.
\]
Thus $p=(0,0,h/D_h)^{\mathsf T}$. The LP passes with $R=\{3\}$,
supplier 1 before supplier 2, and zero demander-only activity.
For uniform reserves, $\kappa=(1-2h)/(4h)$,
$\sigma=D_h/h$, and $1+1/\kappa=\gamma_h$.
This repeats portions of the planner optimum while amplifying
position capacity; it needs no separate completion transaction
and does not assert the smallest approval count.

\paragraph{The interacting four-account cycle.}
For $\widetilde M_h$ and the primitives in
Section~4.3%<xref:sec:interactingexample>
, account 1 requires
$\sum_jb_j\geq a$. Hence $W\leq9a/10\leq9/10$.
The feasible allocation $(1,1/3,1/3,1/3)$ attains equality.
Any maximizer must have $a=1$ and $\sum_jb_j=1$;
the caps $b_j\leq1/3$ then prove uniqueness.

At zero balances, a supplier-1 pair requires both $x\leq y$
and $x\geq2y$, so $x=y=0$. A supplier-2 pair consumes $y$
units of account 3, and a supplier-3 pair consumes $y$ units of
account 4; either therefore requires $y=0$. With $y=0$, any
$x>0$ consumes account 1. Supplier-only moves are not voluntary
because costs are positive. The start is consequently blocked
for every $h\geq0$.

Let $w^j=e_a+e_{b_j}$, each with surplus $9/10$. Their effects are
\[
 \begin{aligned}
 \widetilde M_hw^1&=(0,-1,1,0)^{\mathsf T},&
 \widetilde M_hw^2&=(0,0,-1,1)^{\mathsf T},\\
 \widetilde M_hw^3&=(0,1+h,0,-1)^{\mathsf T}.&&
 \end{aligned}
\]
The optimum decomposes as $u^j=w^j/3$. Only account 2 has a
negative prefix, of size $1/3$; accounts 3 and 4 finance successive
steps. The net release is $p=(0,h/3,0,0)^{\mathsf T}$.
The LP therefore passes with $R=\{2\}$ and supplier order
$(1,2,3)$, giving $\kappa=1/h$, $\sigma=3/h$, and the count
$3\lceil\log(1+h/(3\epsilon))/\log(1+h)\rceil$.

At $h=0$, all columns sum to zero, so $\lambda=\mathbf1$ is a
nonexpansion certificate. Normalize $x+y=1$; voluntary pairs have
$0\leq y\leq10/11$. The negative-resource sums for suppliers
1, 2, and 3 are respectively
\[
 [1-2y]_+ +[3y-1]_+,\qquad |1-2y|+y,
 \qquad [1-2y]_++y.
\]
Their minima are $1/3$, $1/2$, and $1/2$; a normalized
demander-only move has sum one. Hence $\delta_\lambda=1/3$,
$S_\lambda=4$, $\bar v=1$, and $a_C=1/12$.
Since $\widetilde M_h\geq\widetilde M_0$, the general
institutional comparison applies.

At $\rho=1$, $\tau=10^{-6}$, and $F=1/5000$, the conserving
sequential value is bounded above by
$9/10-2\sqrt{10^{-6}(1/12)(9/10)}
=9/10-\sqrt{3/10^7}$.
The joint cost $201/10^6$ is smaller than this cost bound:
squaring compares $40401/10^{12}$ with $300000/10^{12}$.
Thus the exact best joint value $899799/10^6$ is higher.
For $h=1$ and $\epsilon=10^{-6}$,
$2^{18}<1+1/(3\epsilon)<2^{19}$, so $57$ approvals suffice.
The implementation cost is at most $58/10^6$, giving net value
at least $449971/500000$, which exceeds the best joint value by
$143/10^6$. Truncating the last cycle implements the optimum
exactly; the general construction proves all prefix balances,
caps, compensation, and full reserve repayment. These optimized
upper and attainable lower bounds establish the reversal without
asserting exact sequential values.

\section{Optimized implementation losses and institutional thresholds}
\label{app:optimizedloss}

\begin{proof}[Proof of Theorem~6%<xref:thm:optimizedloss>
]
Write $V=V_M>0$. Repayment implies $0<W\leq V$ for any schedule
that can improve on the baseline. All constants below concern the fixed
economy and reserve direction; $\rho>0$ is fixed while $\tau\downarrow0$.

\emph{Nonexpansion.} Fix a strictly positive certificate $\lambda$ and put
$a=\delta_\lambda/(\bar vS_\lambda)>0$, using
\eqref{eq:generalseedconstants}. The blocked-start bound gives
$T\geq aW/\epsilon$, so
\[
 \rho\epsilon+\tau T\geq2\sqrt{\rho\tau aW}.
\]
The function $W-2\sqrt{\rho\tau aW}$ is continuous and convex on
$[0,V]$, with second derivative $\sqrt{\rho\tau a}/(2W^{3/2})$
for $W>0$. Its maximum is therefore attained at an endpoint. Taking the
supremum over allocations and schedules gives
\[
 \mathcal V_M^{\mathrm{seq}}\leq
 \max\{0,V-2\sqrt{\rho\tau aV}\},\qquad
 L_M\geq\min\{V,2\sqrt{\rho\tau aV}\}.
\]
For the upper bound, Theorem~\ref{thm:temporaryreserve} implements a
planner optimum with at most $A/\epsilon+B$ approvals, where
$A=LH^*>0$ and $B=L>0$. Positivity of $H^*$ follows from the blocked
start: otherwise every nonzero decomposition bundle would have no negative
resource effect and would itself be a feasible first move. Choosing
$\epsilon=\sqrt{A\tau/\rho}$ consequently yields
\[
 \min\{V,2\sqrt{\rho\tau aV}\}\leq L_M
 \leq\min\{V,2\sqrt{\rho\tau A}+B\tau\}.
\]
These bounds establish the square-root order and positive adoption for
all sufficiently small $\tau$.

\emph{Executable productivity.} Use the certificate constants
$k,\ell,\alpha,\sigma,\Lambda$ from
Theorem~4%<xref:thm:generalproductive>
. For any strictly positive weights
$\lambda$, the blocked start gives $\delta_\lambda>0$ and the
productive release gives $g_\lambda>0$: otherwise summing its voluntary
bundles would imply $\lambda^{\mathsf T}p\leq0$, contradicting
$p\geq0$, $p\ne0$. Define
\[
 \beta=\log(1+g_\lambda/\delta_\lambda)>0,\qquad
 c=\frac{g_\lambda V}{2\bar vS_\lambda}>0.
\]
An allocation with $W<V/2$ loses at least $V/2$ before costs. For
$W\geq V/2$, the general blocked-start bound instead gives
$T\geq\beta^{-1}\log(1+c/\epsilon)$. With
$y=\beta\rho\epsilon/\tau$ and $t=\beta\rho c/\tau$, use
\[
 y+\log(1+t/y)\geq y-\log y+\log t\geq1+\log t
\]
to bound its implementation cost below. Thus, for
$0<\tau\leq\beta\rho c$,
\[
 L_M\geq\min\left\{\frac V2,
  \frac\tau\beta\left[1+\log\left(\frac{\beta\rho c}{\tau}\right)
                  \right]\right\}.
\]
The inequality covers the baseline and every feasible schedule, hence
also the infimum of their losses; no maximizing schedule is assumed.
The productive construction, applied to a planner optimum with
$\epsilon=\tau/\rho$, supplies the complementary bound
\[
 L_M\leq\min\left\{V,\tau(1+k+\ell)
  +\frac{k\tau}{\log\Lambda}
       \log\left(1+\frac{\alpha\rho}{\sigma\tau}\right)\right\}.
\]
Here $\lceil x\rceil\leq x+1$ bounds the integer approval count.
The two bounds have order $\tau\log(\rho/\tau)$, and the upper
bound tends to zero. This proves the second rate and positive adoption.
\end{proof}

\begin{proof}[Proof of Proposition~7%<xref:prop:feethreshold>
]
Write $S=\mathcal V_M^{\mathrm{seq}}=V_M-L_M$ and
$J=\max\{0,V_M-\tau-F\}$. If $S>0$, direct comparison shows
$J\gtreqless S$ exactly when $F\lesseqgtr L_M-\tau$. If $S=0$,
joint settlement is strictly better exactly when $F<V_M-\tau$;
otherwise both institutions choose the baseline.

A nonzero fully repayable one-approval schedule would be a voluntary
bundle $z$ with $Mz\geq0$, contradicting the blocked start. Therefore
$T\geq2$ for every nonzero schedule and
$S\leq\max\{0,V_M-2\tau\}$. If $S>0$, this gives
$F_M^*=L_M-\tau\geq\tau$. Without a blocked start, every nonzero
schedule still has $T\geq1$. Joint settlement with $F=0$ can implement
its terminal allocation with balanced voluntary transfers, saving
$\rho\epsilon+\tau(T-1)$; the baseline covers nonpositive-surplus
endpoints. Taking suprema proves weak dominance. Under the blocked
start and $V_M>\tau$, the inequality
$V_M-\tau>\max\{0,V_M-2\tau\}\geq S$ proves strict dominance.
\end{proof}

\begin{proof}[Proof of Corollary~8%<xref:cor:optimizedreversal>
]
Positive adoption holds for small $\tau$. Subtracting $\tau$ preserves
both loss orders, so
\[
 0\leq\frac{F_P^*}{F_C^*}
 \leq C\sqrt{\frac\tau\rho}\log(\rho/\tau)\longrightarrow0
\]
for some fixed $C<\infty$. Proposition~7%<xref:prop:feethreshold>
 then
gives the interval and its strict rankings. If $M_P\geq M_C$,
nonnegative activities imply $\mathcal F_C^S(0)\subseteq
\mathcal F_P^S(0)$ and $V_P\geq V_C$, identifying the change as a
productive improvement. Finally, both thresholds tend to zero, so every
fixed positive fee eventually exceeds both.
\end{proof}

\section{Standing balances, reserve borrowing, and institutional choice}
\label{app:standingcapacity}

\subsection{Uniform certificates and primitive verification}

Retain the finite positive-linear economy, complete bilateral access,
positive action caps, balanced compensation, $B\equiv0$, and blocked
start at zero. Let original balances be $sr$, $0\leq s\leq\bar s$,
and rented balances $\epsilon r$, with $r\gg0$. A uniform core and
completion certificate consists of a fixed finite list of voluntary
bundles $u^1,\ldots,u^k$, their sum $d$, a nonnegative resource
release $p=Md$, scalars
$0<\underline\alpha\leq\alpha_s\leq\overline\alpha$, and at most
$\ell$ voluntary completion bundles $v_s^1,\ldots,v_s^{\ell_s}$.
They must sum to a planner optimum
$z_s^*=\alpha_s d+\sum_jv_s^j$ and satisfy
\[
 sr+\alpha_s p+M\sum_{j=1}^t v_s^j\geq0
 \quad(0\leq t\leq\ell_s).
\]
The productive case additionally requires $p\ne0$ and
$M\sum_{j=1}^t u^j\geq-\kappa p$ for every prefix, for a fixed
$\kappa>0$. The completion may consume original balances and preserves all rented capacity.

The certificate has a finite verification on any interval where its
activity bundles are affine. To see this, write the planner objective
as $w^{\mathsf T}z$. The planner linear program and a dual are
\[
 \max\{w^{\mathsf T}z:Mz+sr\geq0,\ 0\leq z\leq\bar z\},
 \qquad
 \min\{s\lambda^{\mathsf T}r+\mu^{\mathsf T}\bar z:
       \mu-M^{\mathsf T}\lambda\geq w,\ \lambda,\mu\geq0\}.
\]
Suppose $\alpha_s$ and a fixed list $v_s^1,\ldots,v_s^\ell$ are
affine, with zero bundles allowed. Fix their transaction types.
Voluntariness, activity caps, nonnegativity, and every completion
inequality are linear in $s$, so checking $s=0$ and $s=\bar s$
checks the entire interval. If one dual-feasible pair $(\lambda,\mu)$
has objective equal to $W(z_s^*)$ at both endpoints, equality holds
throughout the interval. Thus the selected allocation is planner-optimal
for every $s$. The test constructs the certificate from resource effects, preferences,
caps, and linear inequalities.
A stable planner basis often supplies the affine candidate; the
certificate then verifies that its comparative static can be executed.

\subsection{Proof of Theorem~9%<xref:thm:standingcapacity>
}

Put $t=s+\epsilon$. The case $t=0$ has no positive feasible first
approval, and the corresponding envelopes are infinite at
$\epsilon=0$. All statements below concern $t>0$. Planner values
satisfy $0<V(0)\leq V(s)\leq\bar V<\infty$ by monotonicity and
compact action caps.

\emph{Lower bounds uniform across allocations.}
Use all account coordinates and choose strictly positive weights
$\lambda$. A blocked start at zero implies
\[
 \delta=\min_{d'\in\mathcal D}
       \sum_i\lambda_i[-(Md')_i]_+>0,
 \quad S=\lambda^{\mathsf T}r>0,
 \quad g=\max\{0,\max_{d'\in\mathcal D}
                       \lambda^{\mathsf T}Md'\}.
\]
These constants concern the technology at zero balances and remain
fixed when original capacity is introduced. Write each approval as
$x_jd_j$, with $d_j\in\mathcal D$, and let $H_j$ be its available
weighted stock. Feasibility gives
\[
 \delta x_j\leq H_{j-1},\qquad
 H_j\leq H_{j-1}+gx_j,\qquad H_0=tS.
\]
Thus the argument of Theorem~\ref{thm:generalseedbounds} applies
with initial stock $t$, irrespective of how much must eventually
be repaid. In the nonexpanding case take its certificate as
$\lambda$, so $g=0$. Any implemented allocation of surplus $W>0$
then requires
\[
 T\geq\frac{\delta W}{\bar v S t}.
\]
In the productive case $g>0$: otherwise the positive core release
would contradict $\lambda\gg0$. The analogous bound is
\[
 T\geq\frac{1}{\log(1+g/\delta)}
       \log\left(1+\frac{gW}{\bar v S t}\right).
\]

An allocation with $W<V(0)/2$ loses at least $V(0)/2$ in gross
surplus relative to its own planner. For every other allocation,
the bounds yield constants $a_C,a_P,b_P>0$, independent of $s$
and $\epsilon$, such that
\[
 T\geq a_C/t\quad\hbox{or}\quad
 T\geq a_P\log(1+b_P/t),
\]
respectively. The elementary inequality
$\log(1+bx)\geq\min\{1,b\}\log(1+x)$, $b,x>0$, makes the
productive bound comparable to $\log(1+1/t)$. Consequently,
\[
 L_C(s;\rho,\tau)\geq
 \min\{V(0)/2,\min(1,a_C)\Phi_C(s;\rho,\tau)\},
\]
with the analogous inequality for $L_P$ and $\Phi_P$.
Both envelopes converge to zero uniformly in $s$: their maximum
over $s\geq0$ occurs at $s=0$. The fixed first term therefore
ceases to bind for sufficiently small $\tau$, establishing
uniform lower bounds. This step is what permits the institution
to choose an arbitrary allocation.

\emph{Planner implementation under conservation.}
Let
$A_i=\sum_{j=1}^k[-(Mu^j)_i]_+$ and
$A=\overline\alpha\max_i A_i/r_i$.
Choose $n=\max\{1,\lceil A/t\rceil\}$ and repeat the ordered
core list $n$ times, scaling every bundle by $\alpha_s/n$.
The temporary shortfall in each repetition is bounded by
$\alpha_s A_i/n\leq tr_i$. Completed repetitions release
nonnegative capacity. Thus each prefix is feasible and the
completed core leaves $tr+\alpha_sp$.
Execute the completion list. Its prefix balances equal
\[
 \epsilon r+
 \left(sr+\alpha_sp+M\sum_{j=1}^t v_s^j\right)\geq\epsilon r.
\]
The entire rented reserve is therefore returned. The approval
count is at most $k\max\{1,\lceil A/t\rceil\}+\ell$.
All partial activity vectors are bounded by the terminal planner
vector, so the action caps hold throughout.

\emph{Planner implementation with productivity.}
Put $\sigma=\min_{i:p_i>0}r_i/p_i>0$ and
$\Lambda=1+1/\kappa>1$. Since $tr\geq\sigma t p$,
the geometric construction in
Theorem~4%<xref:thm:generalproductive>
 implements the core
$\alpha_s d$ with at most
\[
 k\left\lceil
 \frac{\log(1+\overline\alpha/(\sigma t))}{\log\Lambda}
 \right\rceil
\]
approvals. Truncating the last repetition gives the exact core
quantity. The same completion inequalities then return the rented
reserve. For any fixed upper bound $t\leq t_0<\infty$, these
two constructions give constants $A_C,A_P$ such that their
respective counts are at most $A_C/t$ and
$A_P\log(1+1/t)$; additive constants are absorbed uniformly
on this bounded interval.

The minimizers of the displayed envelopes satisfy
$t_C=\max\{s,\sqrt{\tau/\rho}\}$ and
$t_P=\max\{s,(\sqrt{1+4\tau/\rho}-1)/2\}$.
They are bounded by $\bar s+1$ for small $\tau$.
Evaluating the constructions at these minimizers gives
$L_C\leq\max\{1,A_C\}\Phi_C$ and
$L_P\leq\max\{1,A_P\}\Phi_P$.
If the original definition requires $\epsilon>0$, take an
arbitrarily small positive reserve when the minimizer is zero;
the existing finite schedule is feasible with that extra unused
capacity, and the sequential supremum has the same bound.

\emph{Evaluation of the envelopes.}
For conservation the first-order condition in total stock is
$\rho=\tau/t^2$. The constraint $t\geq s$ gives the formula
in the main text and
$\Phi_C\asymp\tau/(s+\sqrt{\tau/\rho})$.
For productivity it is $\rho=\tau/[t(1+t)]$.
Thus $t_P\sim\tau/\rho$ in the interior, while at $t=s$
the envelope equals $\tau\log(1+1/s)$.
These formulas imply, uniformly on $[0,\bar s]$,
\[
 \Phi_P\asymp
 \tau\log\left(1+\frac1{s+\tau/\rho}\right).
\]
This proves the theorem.

\subsection{Borrowing and institutional choice}

Consider $\eta=\tau/\rho\downarrow0$ and
$\eta\ll s\ll\sqrt\eta$. The productive construction at
$\epsilon=0$ has loss $O(\tau\log(1/s))$, which matches the
optimal order. In a conserving economy, any schedule whose total
loss is $O(\rho\sqrt\eta)$ must have surplus at least $V(0)/2$
for small $\eta$. Its lower bound $T\geq a_C/(s+\epsilon)$
then forces $s+\epsilon=\Omega(\sqrt\eta)$. The reserve cost
simultaneously forces $\epsilon=O(\sqrt\eta)$.
Because $s=o(\sqrt\eta)$, every such schedule has
$\epsilon=\Theta(\sqrt\eta)$. This is a statement about
actual implementations attaining the optimal order, rather than
an identification of the exact reserve optimizer from an envelope.

For Corollary~10%<xref:cor:standingreversal>
, write
$a=\sqrt{\tau/\rho}$. Uniform bounds imply
\[
 \frac{L_P}{L_C}
 \leq C(s+a)\log\left(1+\frac1{s+a^2}\right).
\]
If $s\leq a$, the right-hand side is bounded by
$2Ca\log(1+a^{-2})\to0$. If $s>a$ and $s\to0$, it is
bounded by $2Cs\log(1+s^{-1})\to0$.
This proves the ratio limit along every path $s(\tau)\to0$.

For fixed $s>0$, both envelope minima occur at $\epsilon=0$
for sufficiently small $\tau$. Uniform constants therefore give
\[
 L_C\geq c_C\tau/s,
 \qquad L_P\leq C_P\tau\log(1+1/s).
\]
Since $s\log(1+1/s)\to0$, for all sufficiently small fixed
$s$ the interval
\[
 \left(\max\{0,C_P\log(1+1/s)-1\},\ c_C/s-1\right)
\]
is nonempty. Any $f_s$ in this interval, and the common fee
$F=f_s\tau$, satisfy
$L_P<\tau+F<L_C$. As $V(s)\geq V(0)>0$, both joint
options and the productive sequential option are adopted for
small $\tau$. The comparison follows from the exact joint
value $V(s)-\tau-F$. The fee is common across technologies;
the planner values need not coincide. For a fixed positive fee
independent of $\tau$, both technologies eventually prefer
sequential settlement, because both implementation losses tend
to zero. The corollary identifies the relevant fee interval at
each approval price and its coefficient for fixed standing capacity.

\subsection{Membership of \texorpdfstring{$M_h$}{Mh} and the true planner allocation}

Fix $0\leq h<1/2$ and $0\leq s\leq1/32$, and put
$D_h=2+3h$. Original capacities are $s\mathbf1$.
Weighting the pension and payroll inequalities gives
\[
 b_1+b_2\geq
 \frac{2(1+h)}{D_h}a-\frac{4-2h}{D_h}s.
\]
For example, multiply the pension inequality by $2/D_h$
and the payroll inequality by $2(1-h)/D_h$; the supplier
coefficients are one in both columns, while the demander
coefficient is $-2(1+h)/D_h$.
The resulting employee-surplus bound is
\[
 W(z)\leq
 \left(1-\frac{1+h}{5D_h}\right)a+
 \frac{2-h}{5D_h}s.
\]
Its coefficient on $a$ is positive. The stated $z_{h,s}^*$
attains equality at $a=1$, respects all caps, and has balances
\[
 s\mathbf1+M_hz_{h,s}^*
   =\left(0,0,\frac{h+(6+h)s}{D_h}\right).
\]
This proves planner optimality and the value formula. Positive
dual weights on pension and payroll also imply uniqueness.

Take $d=(1,1/2,1/2)$, whose release is
$p=(h/2,0,0)$, and $\alpha_s=\theta_s=2(1-8s)/D_h$.
The single completion bundle uses supplier 2 and has increments
\[
 x_s=\frac{3h+16s}{D_h},\qquad
 y_s=\frac{2h+(12+2h)s}{D_h}.
\]
It is voluntary because
\[
 x_s-y_s/10
 =\frac{(14/5)h+(74/5-h/5)s}{D_h}>0
\]
whenever $(h,s)\ne(0,0)$; at $(0,0)$ the zero completion is
deleted. Paying $y_s/10$ from the demander to supplier 2
balances compensation and leaves the supplier indifferent.
The proportional construction already supplies voluntary,
balanced compensation at every preceding approval.

The accounts after the proportional stage, before the completion,
are
\[
 (s+\epsilon+h\theta_s/2,\ s+\epsilon,\ s+\epsilon).
\]
The completion consumes $s+h\theta_s/2$ pension units and
$s$ payroll units, and releases $(h+(4-2h)s)/D_h$ position
units. It therefore leaves
\[
 \left(\epsilon,\epsilon,
       \epsilon+\frac{h+(6+h)s}{D_h}\right).
\]
Every account is nonnegative, the reserve is fully repaid, and
all caps hold because adjustments are irreversible and their
terminal vector is $z_{h,s}^*$. Both endpoints are nonnegative, so linear interpolation also preserves
feasibility throughout the atomic approval.

For $h=0$, any fixed voluntary decomposition of $d$ supplies
the divided core construction. For $h>0$, the core construction
in Appendix~\ref{app:generalproductive} supplies the productive
prefix inequality. Its constants depend on $h$ and not on $s$.
The scalar $\theta_s$ is positive and uniformly bounded on
$[0,1/32]$, and the completion is affine in $s$.
Thus both parts of Theorem~9%<xref:thm:standingcapacity>
 apply to
the true planner optima throughout this interval.

\subsection{Exact conserving frontier and institutional cutoffs}
\label{app:standingexact}

\begin{proof}[Proof of Proposition~13%<xref:prop:standingexact>
]
Write $t=s+\epsilon$ and $q=W-H$ for the difference between
payroll and position balances. Conservation gives $|q|\leq3t$
at every intermediate endpoint. Initially $q=0$; the planner
endpoint has $q=-3s$. A supplier-1 increment of size $y$ raises
$q$ by $y/2$, and a supplier-2 increment lowers it by $y/2$.
Hence a schedule with $T$ approvals satisfies
\[
 b_1+b_2=2\sum_{j=1}^T|q_j-q_{j-1}|
       \leq12t(T-1)+6s.
\]
This remains a valid upper bound if demander-only transactions
occur, as they leave $q$ unchanged. At the planner optimum,
$b_1+b_2=1-2s$, proving necessity of the stated reserve.

For attainment put
\[
 r_T=\frac{1-8s}{12(T-1)},\qquad
 e=[r_T-s]_+,\qquad t=s+e=\max\{s,r_T\}.
\]
Choose an alternating supplier order ending with supplier 2.
Set $q_0=0$; at every prefinal step set $q_j=3r_T$ after
supplier 1 and $q_j=-3r_T$ after supplier 2; finally set
$q_T=-3s$. Choose supplier increments
$y_j=2|q_j-q_{j-1}|$. Choose the associated demander increment
to make the pension balance zero at every prefinal endpoint
and equal to $e$ at the last endpoint. At each prefinal endpoint,
payroll and position balances are $(3t+q_j)/2$ and
$(3t-q_j)/2$, respectively. These are nonnegative since
$r_T\leq t$. At completion the balances are $(e,e,e+3s)$.

The supplier quantities sum to
$12r_T(T-1)+6s=1-2s$ and their difference is $-6s$.
Thus their terminal quantities are exactly
$(1/2-4s,1/2+2s)$. The pension identity then gives demander
quantity one. Interior demander increments are $3y_j/4$ for
supplier 1 and $5y_j/4$ for supplier 2. The first increment
adds $t/2$ to the respective expression. The last increment is
$5y_T/4-e/2$. If $r_T\geq s$, the latter equals
$7r_T+8s$ and $y_T=6(r_T+s)$; if $r_T<s$, it equals
$5y_T/4$. All are strictly larger than $y_j/10$.
Paying each supplier $y_j/10$ from its demander therefore
balances compensation and satisfies participation. Every
activity increment is nonnegative; monotonicity and the terminal
planner vector imply all action caps. The final balances return
the full rented reserve. This proves attainment. Solving
$\epsilon_T^*(s)=0$ gives the stated no-rental approval count.
\end{proof}

\paragraph{Optimization over allocations.}
For arbitrary $\rho,\tau>0$, the optimized conserving value is
exactly
\[
 \max\left\{0,\frac{37s}{5}-\tau,
 \sup_{n\geq1}\left[
 \frac9{10}a_n+\frac s5
 -\rho\left[\frac{a_n-8s}{12n}-s\right]_+
 -\tau(n+1)\right]\right\},
\]
where
\[
 a_n=
 \begin{cases}
 1,&\rho\leq54n/5,\\
 \min\{1,(12n+8)s\},&\rho>54n/5.
 \end{cases}
\]
A one-approval allocation uses at most one supplier. Supplier 2
attains the maximum at $(a,b_2)=(8s,6s)$ without rented capacity,
giving the second term in the maximum. A reserve cannot enlarge
the attainable terminal set for one approval because it must
be repaid at that same endpoint.

For $T=n+1\geq2$, the planner dual gives
$W\leq9a/10+s/5$, and total variation gives rented reserve
at least $[(a-8s)/(12n)-s]_+$. The slope of the resulting
net-value bound as a function of $a$ is $9/10$ below its
zero-rental kink and $9/10-\rho/(12n)$ above it. Its maximizer
is $a_n$. Since $a_n\geq8s$, the frontier construction just
proved, with $1$ replaced by $a_n$, attains this bound.
For $a<8s$, the dual gross-surplus bound is smaller than the
feasible one-approval value before costs; such allocations
cannot improve the maximum. At $s=0$ any zero-quantity term
is discarded in favor of the baseline. This proves the formula
over all allocations and schedules.

\begin{proof}[Proof of Corollary~14%<xref:cor:standingcutoffs>
]
At $\rho=12$, the $n=1$ optimum is $a_1=20s$, whereas
every $n\geq2$ chooses $a_n=1$. The one-approval and
two-approval net values are $37s/5-\tau$ and
$91s/5-2\tau$, respectively. Both are strictly below
$9/10+s/5-101/10000$ throughout $s\in[0,1/32]$.
The baseline is also below that positive joint value.
For $n\geq2$, it remains to compare
\[
 C_n(s)=\left[\frac{1-8s}{n}-12s\right]_+
       +\frac{n+1}{10000}
\]
with the joint cost $101/10000$.

At $s_J=1/1208$, the $n=100$ cost equals $101/10000$
with zero rental. Put $A_J=1-8s_J=150/151$.
For every $n\leq99$ rental remains positive and
\[
 C_n(s_J)-\frac{101}{10000}
 =(100-n)\left(\frac{A_J}{100n}-\frac1{10000}\right)>0.
\]
For $n\geq101$, approval costs alone exceed the joint cost.
Since all $C_n$ are nonincreasing, joint settlement strictly
dominates for $s<s_J$.

The $n=99$ cost first equals $101/10000$ at
$s_S=9901/11960000$. For every $n\leq98$,
\[
 C_n(s_S)-C_{99}(s_S)
 =(99-n)\left(\frac{1-8s_S}{99n}-\frac1{10000}\right)>0.
\]
The $n=100$ cost remains exactly $101/10000$ for
$s\geq s_J$, whereas larger counts cost more in approvals
alone. This proves the precise tie interval. For $s>s_S$,
$C_{99}(s)<101/10000$, proving strict sequential dominance.

In $M_{1/4}$ the true planner construction uses proportional
quantity $\theta_s=8(1-8s)/11\leq8/11$. Scale the existing
16-approval proportional construction by $1-8s$ and use its
seed from original balances first, renting any remaining
shortfall. This rents at most $1/144309$. The supplier-2
completion reaches $z_{1/4,s}^*$ in at most 17 approvals.
Its total implementation cost is bounded by
\[
 \frac{17}{10000}+\frac{12}{144309}
 <\frac{101}{10000}.
\]
The productive sequential value therefore strictly exceeds
the joint value at every $s\in[0,1/32]$.
\end{proof}


\subsection{Exact verification of the fixed positive-balance comparison}

Set $s=1/2500$, $\rho=12$, $\tau=1/10000$, and $F=1/100$.
For $h=0$,
\[
 V_{0,s}^*=\frac{11251}{12500},\qquad
 \mathcal V_{0,s}^{\mathrm{joint}}=\frac{44499}{50000}.
\]
The preceding global formula selects $n=100$. Indeed, the untruncated
objective has consecutive differences
$\tau-(1-8s)/[n(n+1)]$, which change sign between $n=99$ and $n=100$;
its reserve remains positive at the minimizer. Zero-rental counts start
at $n=208$ and cost at least $209/10000$. The one- and two-approval
values are smaller. Consequently
\[
 \begin{gathered}
 \epsilon^* = \frac{323}{750000},\qquad T^*=101,\\
 \mathcal V_{0,s}^{\mathrm{seq}}
 =\frac{11251}{12500}-12\frac{323}{750000}-\frac{101}{10000}
 =\frac{221203}{250000},\\
 0.884812<\frac{44499}{50000}=0.889980.
 \end{gathered}
\]

For $h=1/4$,
\[
 z_{1/4,s}^*
 =\left(1,\frac{2492}{6875},\frac{7509}{13750}\right),
 \quad V_{1/4,s}^*=\frac{125007}{137500},
 \quad\mathcal V_{1/4,s}^{\mathrm{joint}}
 =\frac{494473}{550000}.
\]
The proportional target has $\theta_s=4984/6875$.
At $T=9$, the exact frontier gives
\[
 \epsilon_9(1/4,\theta_s)
 =\frac{(1/4)(4984/6875)}{6(3^4-1)}
 =\frac{623}{1650000}<\frac1{2500}=s.
\]
Use that nine-approval construction entirely from original capacity,
leaving the difference $s-623/1650000$ unused in each account.
The completion has
\[
 (x_s,y_s)=\left(\frac{1891}{6875},\frac{101}{550}\right),
 \qquad\hbox{compensation}=\frac{101}{5500}.
\]
Its demander surplus is positive and its supplier surplus is zero.
No reserve is rented, and the ten approvals give
\[
 \mathcal V_{1/4,s}^{\mathrm{seq}}
 \geq\frac{125007}{137500}-\frac{10}{10000}
 =\frac{249739}{275000}=0.9081418181\ldots
 >\frac{494473}{550000}=0.8990418181\ldots.
\]
The construction reaches the new planner optimum exactly. The
conserving value is optimized over all admissible allocations,
reserves, and schedules, so these inequalities prove
the reversal at fixed positive original balances and fixed prices.

\section{Exact verification of the endogenous-allocation reversal}
\label{app:endogenousproof}

This section verifies the productive witness in
Theorem~12%<xref:thm:endogenousreversal>
 using exact arithmetic. Put
$h=1/4$, $\theta=8/11$, and $e=1/144309$. The exact frontier gives
\[
 \gamma_{1/4}=3,\qquad G_{16}(1/4)=2\cdot3^7=4374,
 \qquad
 \epsilon_{16}(1/4,8/11)
 =\frac{(1/4)(8/11)}{6(4374-1)}=e.
\]
The reserve is $e$ in each account. As throughout the paper, the price
$\rho$ applies to the scalar reserve intensity $e$, not separately to
each coordinate; at $\rho=12$ its cost is $12e$.

\subsection{The first sixteen approvals}

Write $(x_t,y_t)$ for the demander and participating supplier increments.
Starting from zero activity and balances $(e,e,e)$, use the following list:
\[
\begin{array}{c|c|c|c}
 t & \text{supplier} & x_t/e & y_t/e\\ \hline
 1 & 1 & 11 & 12\\
 2m,\quad m=1,\ldots,7 & 2 & 30\cdot3^{m-1} &24\cdot3^{m-1}\\
 2m+1,\quad m=1,\ldots,7 &1&42\cdot3^{m-1}&48\cdot3^{m-1}\\
 16&2&26245&26244
\end{array}
\]
The rows indexed by $m$ are interleaved in increasing approval order.
For supplier 1 a pair has resource effect
$(-2x+7y/4,x-y/2,x-y)$; for supplier 2 it has effect
$(-2x+5y/2,x-3y/2,x-y)$. Substitution gives the successive balances
\[
\begin{aligned}
 q^1&=(0,6e,0),\\
 q^{2m}&=(0,0,6\cdot3^{m-1}e), &&m=1,\ldots,7,\\
 q^{2m+1}&=(0,18\cdot3^{m-1}e,0), &&m=1,\ldots,7,\\
 q^{16}&=(13120e,e,e)=(e+1/11,e,e).
\end{aligned}
\]
Every coordinate is nonnegative. The sum of the first sixteen increments is
\[
 z^{16}=(104952e,52476e,52476e)
       =\left(\frac8{11},\frac4{11},\frac4{11}\right).
\]
All increments are nonnegative, so no earlier activity exceeds these totals;
the unit caps therefore hold. These calculations also verify directly that
this schedule realizes the exact-frontier first stage.

\subsection{Completion at the planner optimum and repayment}

Approval 17 uses supplier 2 and increments
\[
 (x_{17},y_{17})=(3/11,2/11),\qquad
 M_{1/4}(3/11,0,2/11)^{\mathsf T}=(-1/11,0,1/11)^{\mathsf T}.
\]
It consumes exactly the pension capacity released by the proportional target.
The final activity and balances are
\[
 z^{17}=\left(1,\frac4{11},\frac6{11}\right)=z_{1/4}^*,
 \qquad q^{17}=(e,e,e+1/11).
\]
The unit caps hold, and repayment of $e$ in each account leaves
$(0,0,1/11)$, precisely the planner's residual ledger. Linear resource
effects make every balance within an approval a convex combination of its
endpoint balances, so intermediate feasibility holds even if an approval is
interpreted as continuous implementation.

At every approval, pay the participating supplier $y_t/10$ and charge the
demander the same amount. The compensation vector sums to zero. The supplier
receives its exact incremental cost, while the demander's gain is
$x_t-y_t/10$. For the four rows of the first-stage list these gains, divided
by $e$, are respectively
\[
 49/5,\qquad (138/5)3^{m-1},\qquad
 (186/5)3^{m-1},\qquad 118103/5.
\]
All are strictly positive. Approval 17 gives the demander $14/55>0$.
Thus every step is voluntary with enforceable balanced compensation; no public
subsidy or decentralized bargaining process is used.

\subsection{Merging the final approvals}

Approvals 16 and 17 involve the same pair. Replacing them by their summed
increment preserves participation, balanced compensation, the terminal
allocation, and nonnegative intermediate balances. It therefore gives a
16-approval witness at the same reserve, with net value
\[
 \frac{436442249}{481030000}+\frac1{10000}
 =\frac{27280647}{30064375}\simeq0.907408.
\]
The main table uses this stronger 16-approval bound. The preceding 17-stage
list separates the proportional-target construction from planner completion.
Rounded decimals are descriptive; all inequalities use exact fractions.

\section{Commitment costs that depend on the allocation}
\label{sec:allocationfees}

The authorization technology in the main text permits certification coverage
to vary with the allocation. The resulting charge $F(z)$ prices the
additional resources used by that joint program, at common administrative
prices across technologies.

A constant fee is not essential to the reversal. Let $F(z)$ be any common
nonnegative fee rule for nonzero allocations, used without change across the
two technologies. Write
\[
 \mathcal V_h^J[F]=\max\left\{0,
    \sup_{z\in\mathcal F^S_h(0)\setminus\{0\}}
          [W(z)-\tau-F(z)]\right\},
 \qquad
 \mathcal F^S_h(0)=\{z\in X:M_hz\geq0\}.
\]
The supremum accommodates fee rules for which a maximizing allocation need
not exist. Sequential compensation and implementation prices are unchanged.

\begin{proposition}[Reversal with allocation-dependent commitment costs]
\label{prop:allocationfees}
In the two economies and at the reserve and approval prices of
Theorem~12%<xref:thm:endogenousreversal>
, suppose every nonzero feasible allocation
has a fee satisfying
\[
 \frac{152309}{96206000}<\underline F
 \leq F(z)\leq\overline F<\frac1{50}.
\]
Then
\[
 \mathcal V_0^J[F]>\mathcal V_0^{\mathrm{seq}},
 \qquad
 \mathcal V_{1/4}^{\mathrm{seq}}>\mathcal V_{1/4}^J[F].
\]
In particular, any common fee rule taking values in $[0.005,0.015]$
preserves the strict reversal. The joint institution need not choose a
frictionless planner optimum.
\end{proposition}

\begin{proof}
In the conserving economy, implementing its planner optimum jointly gives
at least $9/10-\tau-\overline F>8799/10000$, the globally optimized
sequential value. In the productive economy every joint allocation gives
at most $10/11-\tau-\underline F$. The merged 16-approval schedule gives
\[
 \mathcal V_{1/4}^{\mathrm{seq}}
 \geq\frac{10}{11}-\overline C_P,
 \qquad
 \overline C_P=\frac{12}{144309}+\frac{16}{10000}
              =\frac{50603}{30064375}.
\]
The lower fee bound is precisely
$\underline F>\overline C_P-\tau$. Thus this positive sequential value
exceeds both the best nonzero joint value and the zero-cost baseline.
\end{proof}

For an example in which the joint allocation changes, let
\[
 F(z)=\frac1{200}+\left[a-\frac{99}{100}\right]_+,
 \qquad [x]_+=\max\{x,0\},\quad z\ne0.
\]
This continuous nondecreasing fee belongs to the stated interval. For each
$a$, the joint institution uses the cost-minimizing supplier mix in
Proposition~1%<xref:prop:productiveplanner>
. Its objective is
$V_h^*a-\tau-F(z)$, increasing up to $a=99/100$ and decreasing thereafter,
since $0<V_h^*<1$. Consequently its unique allocation is
$(99/100)z_h^*$ in each economy. Joint net values are $0.885900$ at $h=0$
and $0.894900$ at $h=1/4$, and the reversal persists. The administrative coverage charge changes the selected quantity,
while a common resource-price schedule preserves the ranking across technologies.



\section{Allocation distortion and binding participation}
\label{app:allocationdistortion}

\label{sec:economicmargins}

Settlement costs can change the allocation itself, and employee
participation can tighten the physical reserve frontier. The following
examples separate these economic margins. They retain $M_0$, the
$B\equiv0$ mandate, complete information, divisibility, irreversibility,
and employee-balanced compensation. Only preferences and implementation
prices are varied explicitly.

\begin{proposition}[Allocation distortion and a participation frontier]
\label{prop:allocationdistortion}
Let initial balances be zero, let all action caps be one, and require
full reserve repayment in the conserving economy $M_0$.
\begin{enumerate}[label=(\roman*)]
\item Suppose
\[
 U(a)=a-\frac25a^2,\qquad C_j(b_j)=\frac{b_j}{10},\qquad
 \rho=\frac{36}{5},\quad\tau=\frac1{20},\quad
 F=\frac1{100},\quad \mathcal K=0.
\]
The gross planner optimum and the joint-settlement optimum are
$z_1=(1,1/2,1/2)$. The unique sequentially optimal allocation is instead
$z_{7/8}=(7/8,7/16,7/16)$, implemented with four approvals and reserve
$7/288$ in each account. Optimized net surplus is $11/25$ under joint
settlement and $17/160$ under sequential settlement.

\item Suppose instead that $U(a)=a$ and $C_j(b_j)=cb_j$, where
$0<c<1$. For the target $z_1$, the least reserve permitting voluntary
implementation with at most two approvals is
\[
 \epsilon^{\mathrm{vol}}_2(c)
 =\max\left\{\frac1{12},c-\frac34\right\}.
\]
The physical lower bound $1/12$ is attainable voluntarily if and only
if $c\leq5/6$. For $c>5/6$, participation binds and raises the reserve
requirement despite positive gross target surplus $1-c$.
\end{enumerate}
\end{proposition}


\begin{proof}[Proof of Proposition~\ref{prop:allocationdistortion}]
Since every column of $M_0$ sums to zero, terminal feasibility after
repayment forces $M_0z=0$. Thus all feasible allocations are
$z_\theta=(\theta,\theta/2,\theta/2)$, $0\leq\theta\leq1$.

\emph{Part (i).}
Gross surplus is $W(z_\theta)=9\theta/10-2\theta^2/5$, with
derivative $9/10-4\theta/5\geq1/10$. The gross planner
therefore uniquely chooses $\theta=1$. Joint settlement charges
the constant $\tau+F=3/50$ for a nonzero allocation, so it also
chooses $\theta=1$, yielding $1/2-3/50=11/25>0$.
The physical frontier in Proposition~2%<xref:prop:exactproductive>

requires $e\geq\theta/[12(T-1)]$. Its attaining schedule has
ratios $x/y=5/6$ initially, $5/4$ or $3/4$ internally, and
$2/3$ or $7/6$ finally. Hence $x/y\geq2/3$. Because
$U'(a)=1-4a/5\geq1/5$, every approval satisfies
$\Delta U\geq x/5\geq2y/15>y/10=\Delta C$.
Paying the supplier $y/10$ from the demander verifies balanced
compensation and participation, so the physical frontier remains exact.

For $n=T-1\geq1$, net surplus at that frontier is
$f_n(\theta)=(9/10-3/(5n))\theta-2\theta^2/5-(n+1)/20$.
Its unique positive maximizer is
$\theta_n=\min\{1,9/8-3/(4n)\}$. The optimized value,
extended to real $n\geq1$, is
\[
 H(n)=
 \begin{cases}
 \displaystyle\frac58\left(\frac9{10}-\frac3{5n}\right)^2
                    -\frac{n+1}{20},&1\leq n\leq6,\\[4pt]
 \displaystyle\frac12-\frac3{5n}-\frac{n+1}{20},&n\geq6.
 \end{cases}
\]
The branches and first derivatives agree at $6$.
Moreover, $H''(n)=27(1-n)/(20n^4)\leq0$ for $1<n<6$,
and $H''(n)=-6/(5n^3)<0$ for $n>6$.
Thus $H$ is concave, strictly so above $1$. Its values
$H(2)=3/40$, $H(3)=17/160$, and $H(4)=13/128$ satisfy
$H(3)-H(2)=1/32>0$ and $H(4)-H(3)=-3/640<0$.
Nonincreasing discrete differences establish the unique integer
optimum $n=3$; its positive value exceeds the baseline.
Consequently $\theta=7/8$, $T=4$, and
$e=(7/8)/(12\cdot3)=7/288$. Gross surplus $77/160$ less
reserve cost $7/40$ and approval cost $1/5$ gives $17/160$.
The attaining schedule preserves caps, irreversibility, intermediate
feasibility, and repayment; compensation was verified above.

\emph{Part (ii).}
Two approvals must each use one of the suppliers, with supplier
increment $1/2$; there is no spare demander-only approval.
If $x$ is the first demander increment, balanced compensation at
both approvals requires $x\geq c/2$ and $1-x\geq c/2$.
The first account vector is
\[
 \begin{cases}
 (e-2x+3/4,\ e+x-1/4,\ e+x-1/2),&\text{supplier 1 first},\\
 (e-2x+5/4,\ e+x-3/4,\ e+x-1/2),&\text{supplier 2 first}.
 \end{cases}
\]
In the first order, pension feasibility gives
$e\geq2x-3/4\geq c-3/4$. In the second, payroll feasibility
gives $e\geq3/4-x\geq c/2-1/4\geq c-3/4$, using
$c\leq1$. The physical bound also requires $e\geq1/12$.
Every voluntary schedule therefore needs
$e\geq\max\{1/12,c-3/4\}$.

To attain this bound, use supplier 1 first with
$x=\max\{5/12,c/2\}$ and $e=\max\{1/12,c-3/4\}$.
For $c\leq5/6$, the first balances are $(0,1/4,0)$ and
the second demander increment is $7/12$; both participation
inequalities hold. For $c>5/6$, the first balances are
$(0,3c/2-1,3c/2-5/4)\geq0$, and the second demander
increment is $1-c/2\geq c/2$. Pay each supplier $c/2$
from the demander. In this latter case the first demander gain
is zero and the second is $1-c>0$. Both increments respect
the caps; linear account effects preserve feasibility between
endpoints. Since $M_0z_1=0$, final balances are $(e,e,e)$,
permitting full repayment.

At $c=9/10$, the attaining reserve is $3/20$ and the
transaction vectors are $(9/20,1/2,0)$ and $(11/20,0,1/2)$.
Balances move through $(0,7/20,1/10)$ to
$(3/20,3/20,3/20)$. Thus $1/12$ is insufficient despite
positive target surplus $1/10$.
\end{proof}

\section{A fixed-grid reversal with positive original balances}
\label{app:gridrobustreversal}

\begin{proof}[Proof of Corollary~15%<xref:cor:gridrobustreversal>
]
\emph{Productive grid implementation.}
Set $\delta=1/500$ and $e=\delta/4$. Starting from $e\mathbf1$,
execute eight rounds with integer quantities
$(u_1,\ldots,u_8)=(1,2,4,8,16,32,64,56)$.
Each round uses supplier 1 and then supplier 2, with demander and
supplier increments $(u_j\delta,u_j\delta)$ at each approval.
From balances $(p,e,e)$, the successive endpoints are
\[
 (p-u_j\delta/4,e+u_j\delta/2,e),
 \qquad(p+u_j\delta/4,e,e).
\]
Thus a round is feasible if $u_j\delta\leq4p$. Initially
$p=\delta/4$; full rounds double $p$, and the last round's
$56$ units are below the available $128$. Since $\sum_j u_j=183$,
the sixteen approvals reach
$\widehat z_\delta=(366/500,183/500,183/500)$ and balances
$(e+183/2000,e,e)$. A final supplier-2 approval has
$(\Delta a,\Delta b_2)=(134/500,89/500)$ and resource effects
$(-91/1000,1/1000,9/100)$. Hence
\[
 z_\delta=(1,183/500,272/500),\quad
 M_{1/4}z_\delta=(1/2000,1/1000,9/100)\geq0.
\]
All prefixes respect the caps because activities increase to
$z_\delta\leq\mathbf1$; linear interpolation preserves account
feasibility. Paying each supplier its incremental cost and charging
the demander the same amount gives exact balance. Demander gains are
$9u_j\delta/10$ in the rounds and $1251/5000$ at the final approval,
so all trades are voluntary. Terminal balances are
$e\mathbf1+M_{1/4}z_\delta$, permitting full repayment. Adding any
$R^0\geq0$ preserves every constraint. Gross surplus is $909/1000$,
and $12e+17/10000=77/10000$, giving net value $9013/10000$.

\emph{Global conserving bound.}
Relax the grid and start from original balances $s\mathbf1$, with
rented reserve $\epsilon\mathbf1$. For any $T$-approval schedule in
$M_0$, write $a$ and $b=b_1+b_2$ for terminal activities and
$q_t=W_t-H_t$ for the difference between payroll and position balances.
Conservation gives $P_t+W_t+H_t=3(s+\epsilon)$, so
$|q_t|\leq3(s+\epsilon)$ and $q_0=0$. Repayment leaves nonnegative
balances summing to $3s$, hence $|q_T|\leq3s$. Each supplier
increment changes $q_t$ by half its quantity, with sign determined
by supplier identity. For $T\geq2$,
\[
 b=2\sum_{t=1}^T|q_t-q_{t-1}|
 \leq12(s+\epsilon)(T-1)+6s.
\]
The terminal position balance gives $a\leq b+2s$ and
$W(z)\leq9a/10+s/5$. Putting $n=T-1$ therefore yields
\[
 W(z)-12\epsilon-\frac{T}{10000}
 \leq a\left(\frac9{10}-\frac1n\right)
       -\frac{n+1}{10000}
       +\left(\frac{61}{5}+\frac8n\right)s.
\]
The bound remains valid if the implied lower bound on $\epsilon$
is negative. For $n=1$, maximizing the first two terms over
$a\in[0,1]$ gives $-2/10000$. For $n\geq2$, their maximum is
$8799/10000$, since $1/n+(n+1)/10000$ is minimized at $n=100$
with value $201/10000$. The coefficient on $s$ is at most $101/5$.
For $T=1$, $b=2|q_T|\leq6s$ and $a\leq8s$, so net value is at
most $8s-1/10000$, below the same bound. Including the baseline,
\[
 \mathcal V^{\mathrm{seq}}_0(s\mathbf1)
 \leq\frac{8799}{10000}+\frac{101}{5}s.
\]
This also bounds grid settlement; at $\bar s=1/2500$ its right
side is $44399/50000$.

\emph{Joint values and unequal original balances.}
Conserving gross surplus is at most $9/10+s/5$. At $h=1/4$, weight
the pension and payroll constraints by $8/11$ and $6/11$ to obtain
\[
 b_1+b_2\geq\frac{10}{11}a-\frac{14}{11}s,
 \qquad W(z)\leq\frac{10}{11}+\frac7{55}s.
\]
Thus productive joint value, even with divisible activities, is at
most $98889/110000+(7/55)s$. On the grid, gross surplus is a
multiple of $1/5000$, and
\[
 \frac9{10}+\frac{\bar s}{5}<\frac9{10}+\frac1{5000},
 \qquad
 \frac{10}{11}+\frac7{55}\bar s
 <\frac{909}{1000}+\frac1{5000}.
\]
The grid planner values are consequently bounded by $9/10$ and
$909/1000$. The zero-ledger allocations $(1,1/2,1/2)$ and
$z_\delta$ attain these bounds. Subtracting the common positive-package
cost $101/10000$ gives exact joint values $8899/10000$ and
$8989/10000$, both above the baseline.

Attainable values are nondecreasing in each original balance:
adding a nonnegative vector to all ledgers preserves feasibility,
participation, costs, and repayment. For every
$R^0\in[0,\bar s]^3$, zero-ledger witnesses remain available and
$\bar s\mathbf1$ supplies the upper bounds. The conserving and
productive strict margins are respectively
$8899/10000-44399/50000=6/3125$ and
$9013/10000-8989/10000=3/1250$. This proves the corollary without
identifying either optimized sequential value.
\end{proof}

\section{Valued public balances: financing and recovery}
\label{sec:implementation}\label{app:valuedbalances}

The main text retains $B\equiv0$. Here $q=R+Mz$ denotes terminal capacity
after reserve repayment, and $B:\R_+^K\to\R$ is concave, continuously
differentiable, and coordinatewise nondecreasing; derivatives at the boundary
are defined by an extension to a neighborhood. The social benchmark is
\begin{equation}
 \mathcal V_B(R)=\max_{z\in\FS(R)}
 \{W(z)+B(R+Mz)-B(R)\}.
 \tag{S$_B$}\label{eq:socialbenchmark}
\end{equation}
Compactness gives an optimum $z_B^*$, and the baseline gives nonnegative
value. Implementation expenditures still come from separately authorized
operating resources and are deducted from this objective.

\begin{proposition}[Financing and participation]
\label{prop:socialfinance}
A feasible allocation $z$ can make every employee and the public authority
weakly better off precisely when its authorized net public payment $s$ to
employees satisfies
\begin{equation}
 -W(z)\leq s\leq B(R+Mz)-B(R).
 \label{eq:publicfinanceinterval}
\end{equation}
Thus unrestricted financing implements a social optimum voluntarily.
Employee-balanced compensation requires $W(z)\geq0$; public participation
additionally requires $B(R+Mz)\geq B(R)$.
\end{proposition}
\begin{proof}
Employee participation requires $W(z)+s\geq0$ and suffices: for
$\alpha_h\geq0$ summing to one, set
$t_h=\alpha_h[W(z)+s]-f_h(z_h)$. Transfers sum to $s$ and employee
payoffs are nonnegative. Public participation gives the upper bound.
The interval is nonempty whenever social gain is nonnegative.
\end{proof}

The interval concerns compensation before implementation costs. If public
participation must separately cover expenditure $D$, subtract $D$ from its
upper endpoint. Resource capacity does not itself authorize numeraire
payments. For example, take $U(a)=a/2$, $C(b)=b$, unit caps,
$R=(0,1,1)$, and all fiscal coefficients in \textup{(2)}%<xeqref:eq:threeaccounts>
 equal
to one. With $B(q)=2q_P$, social gain is $b-3a/2$ on $b\geq a$.
Its unique optimum $(0,1)$ has employee surplus $-1$ and public gain $2$:
it requires a subsidy $s\in[1,2]$ and cannot be employee-balanced.

\begin{proposition}[Unrestricted competitive implementation]
\label{prop:publiccompetitive}
Every social optimum has prices $p\geq0$ at which employee $h$ maximizes
$f_h(z_h)+p^{\mathsf T}m_hz_h$ on $X_h$, the authority maximizes
$B(q)-p^{\mathsf T}q$ on $q\geq0$, and $q=R+Mz$. Conversely, every
such equilibrium is socially optimal. Payments
\begin{equation}
 t_h=p^{\mathsf T}m_hz_h,\qquad
 t_0=p^{\mathsf T}(R-q)
 \label{eq:publiccompetitivepayments}
\end{equation}
balance across employees and the authority and satisfy participation.
When $B\equiv0$, fixed nonnegative employee rebates totaling
$p^{\mathsf T}R$ make compensation balance among employees.
\end{proposition}
\begin{proof}
At an optimum put $q^*=R+Mz^*$ and $\beta=\nabla B(q^*)\geq0$.
The polyhedral normal-cone condition supplies $\lambda\geq0$ with
$\lambda^{\mathsf T}q^*=0$ and
\[
 \nabla W(z^*)+M^{\mathsf T}(\beta+\lambda)\in N_X(z^*).
\]
Set $p=\beta+\lambda$. Product structure and concavity give every
employee's maximizing condition. For $q\geq0$, concavity also gives
\[
 B(q)-p^{\mathsf T}q-[B(q^*)-p^{\mathsf T}q^*]
 \leq-\lambda^{\mathsf T}q\leq0.
\]
Thus the authority maximizes and markets clear. No strict feasibility is
needed for the polyhedral normal-cone representation. Conversely, summing
the maximizing inequalities against any feasible allocation cancels all
price terms and proves social optimality. Employees can choose zero and
the authority can choose $q=R$, proving participation; payments balance
by market clearing. If $B\equiv0$, then $p^{\mathsf T}q^*=0$, so
employee payments sum to $-p^{\mathsf T}R$. Rebating this nonnegative
receipt preserves choices and participation and gives employee balance.
\end{proof}

These are supporting-price benchmarks with joint effectiveness, not
decentralized bargaining. Adding $B$ does not alter any given target's
resource frontier under unchanged employee participation. It can alter
the chosen target and invalidate an employee-voluntary decomposition:
social optimality gives only
$\nabla W(z_B^*)^{\mathsf T}z_B^*
+\nabla B(R+Mz_B^*)^{\mathsf T}Mz_B^*\geq0$.

\begin{proposition}[Recovery with employee financing]
\label{prop:valuedrecovery}
Retain the access and preference assumptions of
Theorem~\ref{thm:temporaryreserve}. If a social optimum $z_B^*=(a^*,b^*)$
satisfies
\[
 \Gamma(z_B^*)=\sum_iU_i'(a_i^*)a_i^*
                 -\sum_jC_j'(b_j^*)b_j^*\geq0,
\]
then every reserve $\epsilon r\gg0$ permits exact employee-voluntary
bilateral implementation and full repayment. A decomposition into
$L\leq|D||S|+|D|$ bundles $u^\ell$ gives an approval bound
\[
 L\max\{1,\lceil H_B/\epsilon\rceil\},\qquad
 H_B=\max_k\frac{\sum_{\ell=1}^L[-(Mu^\ell)_k]_+}{r_k}.
\]
If public participation is imposed, also require
$B(R+Mz_B^*)\geq B(R)$.
\end{proposition}
\begin{proof}
Put $g_i=U_i'(a_i^*)\geq0$, $c_j^*=C_j'(b_j^*)>0$,
$Q=\sum_i g_i a_i^*$, and $C=\sum_jc_j^*b_j^*$. If $C>0$, then
$Q\geq C$ and the pair bundles
\[
 x_{ij}=a_i^*\frac{c_j^*b_j^*}{C},\qquad
 y_{ij}=b_j^*\frac{g_i a_i^*}{Q}
\]
sum to the target and satisfy $g_ix_{ij}\geq c_j^*y_{ij}$.
A zero supplier increment is a demander-only bundle. If $C=0$, strict
positivity of supplier marginal costs gives $b^*=0$; use demander-only
bundles. Omit zero bundles. Execute the list divided by
$n=\max\{1,\lceil H_B/\epsilon\rceil\}$ in $n$ rounds.
Round-start ledgers lie on the feasible segment from $R$ to $R+Mz_B^*$,
plus the reserve. Within-round decreases are at most
$\sum_\ell[-(Mu^\ell)_k]_+/n\leq\epsilon r_k$, so balances stay
nonnegative and activities never exceed their target quantities.
Concavity and convexity make each benefit increment at least
$g_ix_{ij}/n$ and its paired cost at most $c_j^*y_{ij}/n$.
Reimbursing that cost gives balanced voluntary compensation.
The terminal ledger is $R+Mz_B^*+\epsilon r$, permitting repayment.
\end{proof}

The condition is sufficient, not necessary. A terminal subsidy satisfying
\eqref{eq:publicfinanceinterval} alone does not specify financing for
intermediate approvals. Such a subsidy requires authorized payment timing
and any necessary numeraire bridge finance. Likewise, public participation
at every approval adds $\Delta B_t\geq\Delta s_t$; only terminal public
participation is considered here.

\begin{example}[Public value changes the planner allocation]
\label{ex:valuedproductive}
For $M_{1/4}$ and $B_\lambda(q)=\lambda q_P$, $\lambda\geq0$, the
social optimum is $(1,4/11,6/11)$ when $\lambda<4/55$ and
$(1,1/2,1/2)$ when $\lambda>4/55$; at equality their joining segment
is optimal. The value is $\max\{10/11,9/10+\lambda/8\}$. The
institutional reversal at the prices of
Theorem~12%<xref:thm:endogenousreversal>
 survives for every $\lambda\geq0$.
\end{example}
\begin{proof}
At fixed $a>0$, the feasible supplier triangle has vertices
$a(4/11,6/11)$, $a(1/2,1/2)$, and $a(2/3,1/3)$, with social values
$10a/11$, $a(9/10+\lambda/8)$, and $9a/10$, respectively.
Linearity gives the stated mix, the positive maximal coefficient gives
$a=1$, and equality of the first two coefficients gives $\lambda=4/55$.
For the first optimum, use the merged 16-approval witness of
Theorem~12%<xref:thm:endogenousreversal>
, whose cost is
$12/144309+16/10000<101/10000$. For the second, the exact frontier
implements $z_1$ at cost $9371/5466250<101/10000$. Both return the
reserve, preserve employee participation, and cost less than the best
joint implementation of the same social optimum. In $M_0$, every
repayable terminal balance is zero, so both optimized institutional
values remain unchanged. These comparisons prove the reversal.
\end{proof}


\section{Discrete adjustments and finite approval budgets}
\label{app:indivisibility}

This extension concerns the conserving economy $M_0$. It retains
complete information, irreversible actions, balanced enforceable
compensation, $B\equiv0$, and full reserve repayment. Transfers and
reserve capacity remain divisible.

\paragraph{A finite approval budget.}
With reserve $\epsilon\mathbf1$ and at most $N\geq2$ approvals,
repayment restricts terminal allocations to
$z_\theta=(\theta,\theta/2,\theta/2)$. The exact frontier in
Proposition~2%<xref:prop:exactproductive>
 gives
\[
 \theta_N(\epsilon)=\min\{1,12\epsilon(N-1)\},\qquad
 V_N(\epsilon)=\frac9{10}\theta_N(\epsilon).
\]
Here $V_N$ is gross employee surplus. Necessity follows from the
frontier's lower bound, and its attaining schedule proves sufficiency.
The gross welfare loss is
$\frac9{10}\max\{0,1-12\epsilon(N-1)\}$; the planner optimum
requires exactly $1/[12(N-1)]$ reserve per account. One approval
cannot implement a positive repayable allocation, which uses both
suppliers. Small reserves therefore require a sufficiently large
approval budget.

\paragraph{A fixed adjustment unit.}
Require activity increments in $\delta\mathbb Z_+$, $\delta>0$;
an approval may authorize several units. Let $E_N(\theta,\delta)$
be the infimum reserve per account implementing $z_\theta$ in at
most $N$ grid approvals, and let $E_\infty$ allow any finite number.
An infeasible requirement has value $+\infty$. Exact grid implementation
requires $\theta/(2\delta)$ to be a positive integer.

\begin{proposition}[A discrete reserve floor and the divisible limit]
\label{prop:gridreserve}
For a compatible target,
$E_\infty(\theta,\delta)=\delta/2$, attained with
$\theta/\delta$ approvals. For each fixed $N\geq2$,
\[
 E_N(\theta,\delta)\geq
 \max\left\{\frac\delta2,\frac{\theta}{12(N-1)}\right\},
 \qquad
 E_N(\theta,\delta)\leq\frac{\theta}{12(N-1)}+4\delta
\]
where the upper bound holds for all sufficiently fine compatible
grids. Consequently, for $\delta_m=\theta/(2m)$,
$E_N(\theta,\delta_m)\to\theta/[12(N-1)]$.
\end{proposition}

\begin{proof}
Every grid transaction's resource effects are multiples of
$\delta/2$. From reserve below $\delta/2$, a feasible first
transaction could have no negative effect; conservation would force
every effect to zero. No nonzero bilateral transaction has this
property: the position row requires $x=y$, leaving pension effect
$-y/2$ or $y/2$. A demander-only transaction also cannot have zero
effects. Hence the reserve is at least $\delta/2$.
To attain it, alternate supplier-1 and supplier-2 transactions
$(x,y)=(\delta,\delta)$. From $(\delta/2,\delta/2,\delta/2)$,
the first leaves $(0,\delta,\delta/2)$ and the second restores
the initial balances. Repeat $\theta/(2\delta)$ times. Paying
each supplier $\delta/10$ from the demander gives voluntary,
balanced compensation. Activities increase to the target, all
prefixes are feasible, and the reserve is returned. The divisible
frontier supplies the additional finite-$N$ lower bound.

For the upper bound, take the attaining $N$-approval divisible
schedule from Appendix~\ref{app:exactseed}, at reserve
$e_N=\theta/[12(N-1)]$. Round its cumulative activities down:
$\widetilde z_k^t=\delta\lfloor z_k^t/\delta\rfloor$.
Rounding preserves monotonicity, caps, and each approval's participant
support. Initial and terminal allocations are unchanged because the
target is grid-compatible. For
$r^t=\widetilde z^t-z^t\in(-\delta,0]^3$, the positive
coefficients in each row of $M_0$ sum to at most four, so
$M_0r^t\geq-4\delta\mathbf1$. All rounded prefixes are therefore
feasible with reserve $(e_N+4\delta)\mathbf1$; their linear
interpolations are feasible as well. The unchanged terminal target
has $M_0z_\theta=0$, allowing full repayment.

The attaining schedule has minimum transaction surplus
$\eta=\min_t(x_t-y_t/10)>0$. Each rounded increment differs from
its original by less than $\delta$ in each active coordinate.
Its surplus is therefore at least $\eta-11\delta/10>0$ whenever
$\delta<10\eta/11$. Paying each supplier its incremental cost
gives balanced voluntary compensation in at most $N$ approvals.
The upper bound and the stated limit follow.
\end{proof}

\paragraph{Quantity loss and costs.}
For $0<\delta\leq1/2$, the largest grid-feasible repayable quantity
under the unit caps is
$\theta_\delta=2\delta\lfloor1/(2\delta)\rfloor$; its gross
welfare loss satisfies
$0\leq\frac9{10}(1-\theta_\delta)<\frac95\delta$.
Thus the grid-induced quantity loss and minimum seed vanish with
$\delta$. The minimum-seed construction uses
$\theta_\delta/\delta$ approvals and need not minimize cost when
approvals are priced. For a fixed cap, the rounding construction
instead implements a compatible divisible target with at most
$4\rho\delta$ additional reserve cost and no additional approvals.
Coarse units and finite approval budgets therefore impose distinct
constraints on implementation.

\section{Water trading: data, ordering, and account constraints}
\label{app:waterbenchmark}

\subsection{Data and reconciliation}

The application uses the Victorian Water Register's published reports for
the Goulburn to Murray trade openings in October and December 2025
\citep{DEECA2025October,DEECA2025December}. The extraction preserves the
processing priority, submitted volume, decision, and trading opportunity
before and after each decision. Table~\ref{tab:wateraudit} summarizes the
391 eligible applications. Quantities are represented as integer tenths
of a megalitre, so all balance checks and counterfactual comparisons use
exact arithmetic.

\begin{table}[htbp]
\centering
\caption{Reconciliation of the published processing ledgers}
\label{tab:wateraudit}
\begin{tabular}{lrr}
\toprule
 & October & December \\
\midrule
Eligible applications & 203 & 188 \\
Approved applications & 98 & 117 \\
Requested volume (ML) & 70,222.0 & 68,123.1 \\
Opportunity before processing (ML) & 31,215.1 & 37,647.4 \\
External change during processing (ML) & 0.0 & $-6.2$ \\
Approved volume (ML) & 31,209.8 & 37,639.9 \\
Opportunity after processing (ML) & 5.3 & 1.3 \\
\bottomrule
\end{tabular}
\end{table}

October applications were lodged on 15 October and processed on 16 October.
The opening opportunity of $31,523.4$ ML was reduced by $326.3$ ML of
tagged use and increased by $18.0$ ML of back trade before the eligible
batch was processed. Its initial opportunity was therefore $31,215.1$ ML.
The analysis includes the complete eligible submission batch in Table 5
of the report. It excludes the two ineligible applications and the
separately reported $5.3$ ML application lodged later. Every ordering
comparison uses these same 203 requests.

Recorded priorities and receipt identifiers are unique. Receipt identifiers
were assigned in arrival order, which supplies the chronological
counterfactual. The recorded sequence is reproduced exactly: each of the
203 decisions equals the comparison of requested volume with available
opportunity, and every balance transition equals the approved debit.
Rejected requests leave capacity unchanged and processing continues.
Equal volumes in the smallest first comparison are ordered by recorded
priority. This convention selects one allocation without changing the
maximum count.

\subsection{The count benchmark and its proof}

Fix eligible applications $\mathcal I$ with volumes $v_i>0$, initial corridor
opportunity $R\ge0$, and eligibility held constant across orders. Each
application executes in full or is rejected; none releases corridor
capacity. For a permutation $q$, write
\[
 Y_i(q)=\mathbf 1\!\left\{
 v_i\le R-\sum_{j\prec_q i}v_jY_j(q)\right\},
 \qquad A(q)=\sum_{i\in\mathcal I}Y_i(q).
\]
Write $v_{(1)}\le\cdots\le v_{(n)}$, $S_0=0$, and
$S_k=\sum_{j=1}^k v_{(j)}$.

\begin{proposition}[Application count and additional capacity]
\label{prop:watercardinality}
The maximum count under sequential processing is
\[
 k^*(R)=\max\{k\in\{0,\ldots,n\}:S_k\le R\}.
\]
Smallest first attains this count. Joint selection attains the same count,
and every jointly feasible subset is sequentially feasible in every
order. The minimum additional corridor capacity required for $k$
applications is $(S_k-R)_+$. The maximizing order need not be unique.
\end{proposition}

\begin{proof}
Every subset of $k$ applications has volume at least $S_k$, so no feasible
selection contains more than $k^*(R)$ applications. In increasing volume
order, every prefix through $k^*(R)$ fits. The next request, if present,
exceeds residual capacity, as do all subsequent requests. The rule
therefore attains the upper bound. This selected subset is jointly
feasible. Conversely, positive volumes imply that every prefix of a
jointly feasible subset uses at most the subset's total volume. Every
order of that subset is consequently feasible. Additional capacity $c$
permits $k$ requests exactly when $R+c\ge S_k$, giving the stated minimum.
Any permutation of an optimal subset placed first attains the same count,
so uniqueness is unnecessary.
\end{proof}

For October the exact certificate is
\[
 S_{159}=31{,}193.0\le31{,}215.1<31{,}693.0=S_{160}.
\]
Thus the scalar maximum is 159 applications. Executing 160 would require
at least $477.9$ ML of additional corridor capacity. This increment is
consumed trading opportunity, distinct from the repayable temporary
reserve in the main model. Twenty three requests have the cutoff volume
of $500$ ML, reinforcing the nonuniqueness of the maximizing allocation.

Relative to recorded priority, the selected allocation retains 79
approved applications, replaces 19, and admits 80 previously refused
applications. The net increase is 61. Its executed volume is $16.8$ ML
lower. These selection margins distinguish a larger count from a Pareto
improvement. All 203 applications underwent administrative processing;
successful applications are also distinct from the costly approval
steps priced by the main model.

The count objective gives each application unit weight. With gains
$g_i$, the relevant allocation objective instead maximizes
$\sum_i g_i x_i$ subject to $\sum_i v_i x_i\le R$ and
$x_i\in\{0,1\}$. Smallest first solves that problem for equal positive
gains, whereas heterogeneous gains generally require another selection.
The October records classify 191 applications as related party trades
and report a zero declared price for 174. Declared prices therefore do
not identify gains from trade, and the empirical outcomes remain
application counts and executed volumes. Applications also need not
represent distinct traders.

\subsection{A conditional randomization benchmark}

The ACCC independently observed the October randomization procedure
\citep{ACCC2025October}. To characterize ordering variation within the
fixed application menu, the replication draws 100,000 uniform
permutations with NumPy PCG64 and seed 20251015. This stipulated
benchmark yields a mean of $93.66423$ approved applications, a standard
deviation of $9.39806$, and a median of 94. The central $95\%$ interval
runs from 75 to 112. The Monte Carlo standard error of the mean is
$0.02972$ applications.

The observed count of 98 is ordinary within this distribution. The
comparison with smallest first reflects a systematic difference in
selection rules, rather than an unusually unfavorable lottery draw.
The simulated mean gap from 159 is $65.33577$ applications. These
quantities describe a conditional simulation distribution; they are not
estimates of a causal population effect. A permanent policy change can
change submissions and request sizes, including incentives to split
applications. The exercise identifies the ordering margin for the
observed requests, while the institution's equity and allocation
objectives govern the policy interpretation.

\subsection{Account restrictions and the connection to the theory}

December provides an additional ledger check. The opening release of
$38,085.0$ ML was reduced by $437.6$ ML of tagged use before processing,
giving $37,647.4$ ML. A further $6.2$ ML debit occurred between priorities
70 and 71. The application at priority 45 requested $62.3$ ML and was
refused while opportunity remained $23,824.0$ ML. Footnote 8 attributes
the refusal to an insufficient seller account balance
\citep{DEECA2025December}. Incorporating that eligibility restriction
and the recorded external debit reproduces all 188 decisions and balance
transitions, including the $1.3$ ML final opportunity.

A seller's allocation account records private holdings. Its restriction
establishes that aggregate corridor opportunity alone is insufficient
for execution; it does not identify a separate public account of the
main model. Recoverability would require a feasible transfer that
credits the seller before processing its request. The reports do not
provide the complete seller ledgers needed to establish such a path
or to reconstruct eligibility under alternative October orders.
Accordingly, 159 is the exact capacity benchmark conditional on fixed
eligibility and an upper benchmark for a system with additional
restrictions that only exclude transactions.

The application gives the theory a concrete institutional margin:
almost identical aggregate utilization can support markedly different
numbers of executed requests. In this scalar setting, optimal joint
and sequential selection coincide. The main theory adds the economic
content of separate accounts and productive releases, characterizing
the cost of an admissible schedule when the authority chooses its
allocation. The empirical application establishes the importance of
selection order with positive existing capacity and identifies the
account information needed to investigate those further mechanisms.

\section{Relationship to existing work}
\label{sec:literature}\label{sec:literaturecomparison}

Table~\ref{tab:literatureboundary} relates the implementation-loss results to
their mathematical and economic antecedents. The divided-list construction
in \citet[Lemma~13]{FracaHaddad2015} provides an inverse-buffer upper bound.
Our analysis adds a matching nonexpansion lower bound and an exact frontier.
The economic comparison allows each institution to choose its allocation
subject to voluntary participation and balanced compensation. Standing
capacity and productive releases then determine reserve borrowing and the
maximum fee the authority will pay for joint authorization.

\paragraph{The precise continuous-net connection.}
Fix a finite list of voluntary bundles in the linear model. Represent each
bundle $u$ by a transition with
\[
 \operatorname{Pre}(u)=([-Mu]_+,u),\qquad
 \operatorname{Post}(u)=([Mu]_+,0).
\]
The state $(q,\bar z-z)$ records balances and remaining activity caps.
Firing in amount $s\geq0$ is equivalent to $q+sMu\geq0$ and
$z+su\leq\bar z$. Voluntariness and balanced compensation scale with $u$.
Rational coefficients permit transition rescaling to integer arc weights.
This represents a fixed list of atomic approvals, not every continuously
chosen bilateral bundle by one fixed finite net. Splitting an approval into
separate activity moves can change feasibility and the approval count.

\begin{longtable}{@{}p{0.45\linewidth}p{0.50\linewidth}@{}}
\caption{Antecedents and scope of the quantitative results}
\label{tab:literatureboundary}\\
\toprule
Established result & Result and boundary here\\
\midrule
\endfirsthead
\toprule
Established result & Result and boundary here\\
\midrule
\endhead
\emph{Reachability and inverse bounds.} Activation and divided execution
\citep[Lemma~13, Theorem~20]{FracaHaddad2015}. & Voluntary atomic approvals,
repayment, a matching nonexpansion lower bound, and an exact frontier for
$M_h$.\\[3pt]
\emph{Amplification.} Enabled resource-producing directions
\citep[Theorem~25]{FracaHaddad2015}. & A finite executable release and
completion give sharp logarithmic order; this is a sufficient class,
not a classification of all technologies.\\[3pt]
\emph{Costs.} Square-root inventory minimization \citep{Baumol1952}. &
Theorem~6%<xref:thm:optimizedloss>
 bounds losses after allocation choice:
$\Theta(\sqrt\tau)$ versus $\Theta(\tau\log(1/\tau))$ at fixed $\rho$.\\[3pt]
\emph{Settlement welfare and financing.} Credit constraints and intraday
finance \citep{KahnRoberds2001}; gridlock, collateral, and delay
\citep{KahnMcAndrewsRoberds2003}; queue incentives
\citep{MartinMcAndrews2008}; planner comparisons
\citep{AtalayMartinMcAndrews2010}. &
Productive releases change the cost of an authority-selected execution
schedule and its willingness to pay for joint authorization.\\
\bottomrule
\end{longtable}

Theorem~6%<xref:thm:optimizedloss>
 establishes matching orders for optimized
implementation losses at a blocked start. Theorem~9%<xref:thm:standingcapacity>

extends the quantitative comparison to positive standing balances under its
uniform execution conditions. The fee thresholds translate these losses into
institutional choice after allocation optimization. The execution certificate
provides a sufficient class for the productive bound, while
Proposition~2%<xref:prop:exactproductive>
 supplies the exact constants for $M_h$.

\end{appendix}
\begin{thebibliography}{99}
\bibitem[Atalay et~al.(2010)]{AtalayMartinMcAndrews2010}
\textsc{Atalay, Enghin, Antoine Martin, and James McAndrews} (2010):
``The welfare effects of a liquidity-saving mechanism,'' Staff Report 331,
Federal Reserve Bank of New York, June 2008, revised January 2010.
\url{https://www.newyorkfed.org/research/staff_reports/sr331.html}.


\bibitem[ACCC(2025)]{ACCC2025October}
\textsc{Australian Competition and Consumer Commission} (2025):
``Goulburn to Murray trade opening, 15 October 2025: ACCC independent observer
of the execution of randomisation tool,'' statement, 24 October.
\url{https://waterregister.vic.gov.au/images/documents/ACCC_statement_randomisation_october_2025.pdf}.

\bibitem[Baumol(1952)]{Baumol1952}
\textsc{Baumol, William J.} (1952): ``The transactions demand for cash: An inventory theoretic approach,'' \emph{The Quarterly Journal of Economics}, 66(4), 545--556. DOI: \url{10.2307/1882104}.

\bibitem[DEECA(2025a)]{DEECA2025October}
\textsc{Department of Energy, Environment and Climate Action} (2025a):
\emph{Goulburn-to-Murray Post Trade Opening Report: October 2025}.
Victorian Water Register, 23 October.
\url{https://waterregister.vic.gov.au/images/documents/Goulburn-to-Murray-Post-Trade-Opening-Report-October-2025.pdf}.

\bibitem[DEECA(2025b)]{DEECA2025December}
\textsc{Department of Energy, Environment and Climate Action} (2025b):
\emph{Goulburn-to-Murray Post Trade Opening Report: December 2025}.
Victorian Water Register, 19 December.
\url{https://waterregister.vic.gov.au/images/documents/Goulburn-to-Murray-Post-Trade-Opening-Report-December-2025.pdf}.

\bibitem[Fraca and Haddad(2015)]{FracaHaddad2015}
\textsc{Fraca, Estíbaliz, and Serge Haddad} (2015): ``Complexity analysis of continuous Petri nets,'' \emph{Fundamenta Informaticae}, 137(1), 1--28. DOI: \url{10.3233/FI-2015-1168}.

\bibitem[Kahn and Roberds(2001)]{KahnRoberds2001}
\textsc{Kahn, Charles M., and William Roberds} (2001): ``Real-time gross settlement and the costs of immediacy,'' \emph{Journal of Monetary Economics}, 47(2), 299--319. DOI: \url{10.1016/S0304-3932(01)00047-2}.

\bibitem[Kahn et~al.(2003)]{KahnMcAndrewsRoberds2003}
\textsc{Kahn, Charles M., James McAndrews, and William Roberds} (2003): ``Settlement risk under gross and net settlement,'' \emph{Journal of Money, Credit and Banking}, 35(4), 591--608.
\url{https://www.newyorkfed.org/research/staff_reports/sr86.html}.

\bibitem[Martin and McAndrews(2008)]{MartinMcAndrews2008}
\textsc{Martin, Antoine, and James McAndrews} (2008): ``Liquidity-saving mechanisms,'' \emph{Journal of Monetary Economics}, 55(3), 554--567. DOI: \url{10.1016/j.jmoneco.2007.12.011}.

\end{thebibliography}
\end{document}
